% =============================================================================
%  chapter2_body.tex --- reorganized Chapter 2, full body.
%  Content copied/adapted from:
%    sections/RamificationAfterBlowupIntroduction.tex        (opening)
%    sections/Preliminaries/LocalRamificationPreliminaries.tex  (2.1, 2.3.1-2)
%    sections/RamificationAfterBlowupIsTame.tex              (2.2, 2.4, 2.5, 2.6)
%    sections/RamificationAfterBlowupPPower.tex              (2.3.3, 2.5.3)
%    witt/WittVectorPrerequisites.tex                        (\input, unchanged)
%  References prefixed "main:" point outside this chapter and resolve against
%  ../main.aux (xr).
% =============================================================================

\chapter{Ramification After Blowup Is Tame}

% ---------------------------------------------------------------------------
%  Chapter opening: setup, the torsor formulation, roadmap.
%  [Source: RamificationAfterBlowupIntroduction.tex; roadmap rewritten.]
% ---------------------------------------------------------------------------

Throughout this chapter, let $C$ be a smooth, projective, geometrically connected
curve over $k$. Fix a modulus

\[
    \mdls = \sum_i n_i P_i
\]

on $C$, and set $U = C \setminus \mdls$.

For each single-point submodulus $\ndls = n_i P_i$ of $\mdls$, set
$d_\ndls = \deg(\ndls)$. The effective divisor $\ndls$ determines a
$k$-rational closed point

\[
    Z_\ndls \hookrightarrow C^{(d_\ndls)}
        = \operatorname{Sym}^{d_\ndls}_k(C).
\]

Let
\[
    \pi_\ndls \colon
    X_\ndls := \operatorname{Bl}_{Z_\ndls}\bigl(C^{(d_\ndls)}\bigr)
    \longrightarrow C^{(d_\ndls)}
\]

be the blowup at this point. We denote its exceptional divisor by

\[
    E_\ndls := Z_\ndls \times_{C^{(d_\ndls)}} X_\ndls
\]

By
\Cref{main:prop:relative_symmetric_power_smooth,main:prop:exceptional_divisor_smooth_point},
$E_\ndls\cong\mathbb P_k^{d_\ndls-1}$; in particular, it is irreducible and
has codimension $1$ in $X_\ndls$. We denote its unique generic point by
$\eta_\ndls$, so that $\overline{\{\eta_\ndls\}}=E_\ndls$. Thus we have the
Cartesian diagram

\begin{equation}\label{equation:blowup_diagram_ndls}
\begin{tikzcd}
E_\ndls = \overline{\{\eta_\ndls\}}
    \arrow[r, hook] \arrow[d]
  & X_\ndls \arrow[d, "\pi_\ndls"] \\
Z_\ndls \arrow[r, hook]
  & C^{(d_\ndls)}.
\end{tikzcd}
\end{equation}

The purpose of this chapter is to prove
\Cref{main:theorem:SymmetricPowerOfSheavesIsTamelyRamifiedReduction}. Let $\cF$ be
the rank-one $\Lambda$-local system appearing in that theorem, and let

\[
    \cP = \operatorname{Fr}(\cF)
\]

be its $G = \Lambda^\times$-torsor of frames. By
\Cref{main:prop:torsor_module_equivalence,main:definition:symmetric_power_local_system},
there is a canonical identification

\[
    \operatorname{Fr}\bigl(\cF^{[d_\ndls]}\bigr)
        \cong \cP^{[d_\ndls]}.
\]

In particular, the two objects have the same ramification. It therefore
suffices to prove the following torsor formulation.

\begin{theorem}\label{theorem:torsors_after_blowup_is_tame}
Let $G$ be a finite abelian group, and let $\cP \to U$ be a $G$-torsor whose
ramification is bounded by $\mdls$. For every single-point submodulus
$\ndls = n_iP_i$ of $\mdls$, the symmetric-power torsor
$\cP^{[d_\ndls]}$ is tamely ramified at $\eta_\ndls$.
\end{theorem}

The chapter is organized around a single observation: the proof of
\Cref{theorem:torsors_after_blowup_is_tame} is essentially the \emph{same
computation carried out three times}, once for each of the cyclic building
blocks of $G$:
\[
    G = \Z/e\Z \ \ (\gcd(e,p)=1), \qquad
    G = \Z/p\Z, \qquad
    G = \Z/p^r\Z,
\]
governed respectively by Kummer theory, Artin--Schreier theory, and
Artin--Schreier--Witt theory. We therefore proceed as follows.
\Cref{section:ramification_of_discrete_valuations} reviews the ramification
theory of discrete valuations, and
\Cref{section:ramification_of_torsors} defines what it means for a $G$-torsor
to be (tamely) ramified, or to have ramification bounded by a modulus, and
establishes the basic stability properties of these notions.
\Cref{section:cyclic_extension_theories} then presents the three cyclic
theories \emph{in parallel}: each one classifies cyclic extensions by an
explicit equation, and reads the ramification off the valuation of a reduced
representative; a comparison table summarizes the correspondence.
\Cref{section:blowup_prelude} is a prelude collecting the geometric
background propositions shared by the three computations --- the local ring at
the generic point of the exceptional divisor, the formal--local invariance of
the problem, the reduction to $\bP^1$, and the invariant-theoretic
description of the symmetric-power cover. With these tools in place,
\Cref{section:symmetric_power_computations} carries out the three
computations side by side, following the same three-step skeleton
(\emph{realize} the local torsor by an explicit equation over $\bG_m$;
\emph{identify} the symmetric-power cover at $\eta_\ndls$; \emph{count}
valuations to conclude), and then assembles the general case of
\Cref{theorem:torsors_after_blowup_is_tame} by decomposing $G$ into primary
cyclic factors. Finally, \Cref{subsection:ram_prod_analysis} proves auxiliary
results on products of blowups used in \Cref{main:section:gcft}.

% ---------------------------------------------------------------------------
%  2.1  Ramification of Discrete Valuations
%  [Source: LocalRamificationPreliminaries.tex, minus the Kummer/AS theorems
%   (moved to 2.3) and the Witt input (moved to 2.3.3); duplicate
%   Cohen-structure paragraph merged.]
% ---------------------------------------------------------------------------

\section{Ramification of Discrete Valuations}\label{section:ramification_of_discrete_valuations}
This section reviews the ramification of discrete valuations, beginning with the general theory of \textit{discrete valuation rings} and their
integral closures in finite separable extensions.
We then specialize to complete rings in Galois extensions to describe the ramification filtration via lower and upper numbering,
following the treatments in \stackstag{0EXQ} and \cite{serreLF}.

\begin{definition}[\stackstag{09E4}]\label{definition:weakly_unramified_and_tamely_ramified_DVR}
    We say that $A \to B$ or $A \subset B$ is an extension of discrete valuation rings if $A$ and $B$ are discrete valuation rings and $A \to B$ is injective and local. In particular, if $\pi _ A$ and $\pi _ B$ are uniformizers of $A$ and $B$, then $\pi _ A = u \pi _ B^ e$ for some $e \geq 1$ and unit $u$ of $B$. The integer $e$ does not depend on the choice of the uniformizers as it is also the unique integer $\geq 1$ such that
\[ \mathfrak m_ A B = \mathfrak m_ B^ e \]
The integer $e$ is called the \textit{ramification index} of $B$ over $A$. We say that $B$ is \textit{weakly unramified} over $A$ if $e = 1$.
If the extension of residue fields $\kappa _ A = A/\mathfrak m_ A \subset \kappa _ B = B/\mathfrak m_ B$ is finite,
then we set $f = [\kappa _ B : \kappa _ A]$ and we call it the \textit{residual degree} or residue degree of the extension $A \subset B$.

Note that we do not require the extension of fraction fields to be finite.
\end{definition}

Now let $A$ be a discrete valuation ring with fraction field $K$, let  $L/K$ be a finite separable field extension and
let $B \subset L$ be the integral closure of $A$ in $L$.
Then the ring extension $A \subset B$ is finite, hence $B$ is Noetherian.
The dimension of $B$ is $1$, hence $B$ is a Dedekind domain.
Let $\mathfrak m_1, \ldots , \mathfrak m_ n$ be the maximal ideals of $B$ (i.e., the primes lying over $\mathfrak m_ A$).
We obtain extensions of discrete valuation rings

\[ A \subset B_{\mathfrak m_ i} \]
and hence ramification indices $e_ i$ and residue degrees $f_ i$. We have

\[ [L : K] = \sum \nolimits _{i = 1, \ldots , n} e_ i f_ i \]
Note that if $A$ is henselian (e.g. $A$ is complete) then $n=1$.

\begin{definition}[\stackstag{09E9}]\label{definition:tamely_ramified_dvrs}
Let $A$ be a discrete valuation ring with fraction field $K$. Let $L/K$ be a finite separable extension. With $B$ and $\mathfrak m_ i$, $i = 1, \ldots , n$
as above, we say the extension $L/K$ is
\begin{enumerate}
    \item unramified with respect to $A$ if $e_ i = 1$ and the extension $\kappa (\mathfrak m_ i)/\kappa _ A$ is separable for all $i$,
    \item tamely ramified with respect to $A$ if either the characteristic of $\kappa _ A$ is $0$ or the characteristic of $\kappa _ A$ is $p > 0$, the field extensions $\kappa (\mathfrak m_ i)/\kappa _ A$ are separable, and the ramification indices $e_ i$ are prime to $p$, and
    \item totally ramified with respect to $A$ if $n = 1$ and the residue field extension $\kappa (\mathfrak m_1)/\kappa _ A$ is trivial.
\end{enumerate}
If the discrete valuation ring $A$ is clear from context, then we sometimes say $L/K$ is unramified, totally ramified, or tamely ramified for short.
\end{definition}

For the remainder of this section, assume $A$ is a \textit{complete discrete valuation ring} with uniformizer $\pi$ and a perfect residue
field $\kappa$. In the equal characteristic case, where $\text{char}(A) = \text{char}(\kappa) = p > 0$, the
Cohen Structure Theorem implies that $A$ contains a coefficient field $k \cong \kappa$, such that $A \cong k[[\pi]] \cong k[[t]]$.
Consequently, the fraction field is given by $K = k((\pi))$.
Let $L/K$ be a finite Galois extension with Galois group $G$.
The integral closure $B$ of $A$ in $L$ is itself a complete discrete valuation ring.
Since the residue field $\kappa$ is perfect, the extension of residue fields is separable, and there exists a uniformizer
$\Pi \in B$ such that $B = A[\Pi]$.

The \textit{lower numbering ramification filtration} of $G$, denoted by $(G_i)_{i \geq -1}$, is defined as:
\[
G_i = \{ \sigma \in G \mid v_B(\sigma (x) - x) \geq i + 1 \text{ for all } x \in B \}
\]
where $v_B$ is the valuation on $L$ associated with $B$. In this indexing, $G_{-1} = G$ and $G_0$
corresponds to the \textit{inertia group} of the extension $L/K$. Note that $L/K$ is unramified if and only if
$G_0 = \{1\}$, and it is tamely ramified if and only if $G_1 = \{1\}$.

A standard result in the theory of local fields (complete local fields with a discrete valuation and perfect residue field)
ensures that the condition in the definition of $G_i$ need only be checked for the uniformizer $\Pi$ of $B$.
Specifically, if we define the function
\[
i^L_K(\sigma) = v_B(\sigma(\Pi) - \Pi)
\]
for $\sigma \in G$, then $G_i = \{ \sigma \in G \mid i^L_K(\sigma) \geq i + 1 \}$.
The subgroups $G_i$ are normal in $G$ and become trivial for sufficiently large $i$.

Consider a tower of fields $K \subset E \subset L$, and let $H = \text{Gal}(L/E) \subset G$. The lower numbering is compatible with subgroups:
\[
G_i \cap H = H_i \quad \text{for all } i \geq -1
\]
which follows from the identity $i^L_E = i^L_K|_{H}$.
However, the lower numbering does not behave as simply with respect to quotients. If $H$ is normal in $G$, the image of $G_i$ in $G/H$ is generally not $(G/H)_i$.
Instead, we have $G_i H / H = (G/H)_j$, where the index $j$ is given by the formula:
\[
j = \frac{1}{e_{L/E}} \sum_{\tau \in H} \min(i^L_K(\tau), i+1) - 1
\]


In the literature, one reindexes the ramification groups by defining the Herbrand function $\phi_{L/K} : [-1, \infty ) \to [-1, \infty )$:
\[ \phi_{L/K}(i) = \frac{1}{e_{L/K}}\sum_{\sigma \in G} \min(i^L_K(\sigma), i+1)-1 = \int _0^i \frac{1}{[G_0 : G_t]} dt \]

This function is continuous, strictly increasing, and piecewise linear, making it a bijection.
It satisfies the composition law $\phi_{L/K} = \phi_{E/K} \circ \phi_{L/E}$ for a tower of extensions $K \subset E \subset L$.
Using this function, we define the \textit{upper numbering ramification groups} as:
\[
G^i = G_{\phi_{L/K}^{-1}(i)}
\]
The primary advantage of this reindexing is its compatibility with quotients: for any normal subgroup $H \subset G$, we have:
\[
G^i H / H = (G/H)^i
\]
for all $i \geq -1$.

We record, for orientation, the structure of the inertia group and of the
graded pieces of the filtration; these facts are not used in the sequel, but
they explain the dichotomy --- prime-to-$p$ versus $p$-power --- that shapes
the whole chapter.

\begin{prop}[\Cite{serreLF}, IV Corollary 3]
    The quotients $G_i/G_{i+1}$, $i \geq 1$, are abelian groups, and are direct products of cyclic groups of order $p$.
    The group $G_1$ is a $p$-group.
\end{prop}

\begin{prop}[\Cite{serreLF}, IV Corollary 4]
    The inertia group $G_0$ has the following property:
    \textit{It is the semi-direct product of a cyclic group of order prime to $p$ with a normal subgroup whose order is a power of $p$.}
\end{prop}

% ---------------------------------------------------------------------------
%  2.2  Ramification of G-Torsors
%  [Source: RamificationAfterBlowupIsTame.tex lines 1--121, moved BEFORE the
%   Kummer/AS/ASW theories per the reorganization.]
% ---------------------------------------------------------------------------

\section{\texorpdfstring{Ramification of $G$-Torsors}{Ramification of G-Torsors}}\label{section:ramification_of_torsors}

Generally speaking, assume we are given a locally Noetherian scheme $X$,
a dense open $U \subset X$, a finite \'etale morphism $f: Y \to U$ and a prime divisor $Z \subset X$ such that $Z \cap U = \emptyset$
and the local ring $\Oo_{X, \xi}$ of $X$ at the generic point $\xi$ of $Z$ is a discrete valuation ring.
Setting $K_\xi = \Frac(\Oo_{X, \xi})$ we obtain a cartesian square
\[ \xymatrix{ \mathop{\mathrm{Spec}}(K_\xi ) \ar[r] \ar[d] & U \ar[d] \\ \mathop{\mathrm{Spec}}(\mathcal{O}_{X, \xi }) \ar[r] & X } \]
of schemes.
In particular, we see that $Y \times _ U \mathop{\mathrm{Spec}}(K_\xi )$ is the spectrum of a finite separable algebra $L_\xi /K_\xi $.
\begin{definition}\label{definition:tamely_ramified_in_codim_1}
We say:
\begin{enumerate}
    \item $Y$ is \textit{unramified} over $X$ at $(Z, \xi)$ resp. $Y$ is \textit{tamely ramified} over $X$ at $(Z, \xi)$
        if $L_\xi /K_\xi $ is unramified, resp. tamely ramified with respect to $\mathcal{O}_{X, \xi }$
        (\Cref{definition:tamely_ramified_dvrs})

    \item $Y$ is unramified over $X$ in codimension $1$, resp. $Y$ is tamely ramified over $X$ in codimension $1$ if $L_\xi /K_\xi $ is unramified, resp. tamely ramified
    with respect to $\mathcal{O}_{X, \xi }$ for every $(Z, \xi )$ as above.
\end{enumerate}

More precisely, we decompose $L_\xi $ into a product of finite separable field extensions of $K_\xi $ and we require each of these to be unramified, resp. tamely ramified with
respect to $\mathcal{O}_{X, \xi }$.

\end{definition}

Now back to our original assumptions, $C$ is a projective smooth geometrically connected curve
over a field $k$ and $U=C\setminus \mdls$. Every $(Z, \xi)$ is now a closed point $P$ and $\Oo_{C, P}$ is a discrete valuation ring.
Assume $f: Y \to U$ is actually a $G$-torsor $\cP \to U$ in $U_{\text{\'Et}}$ for $G$ a finite abelian group.

By passing to the completion $\widehat{\Oo}_{C,P}$, we obtain the complete local field $\widehat{K}_P$.
Restricting the torsor $\cP$ to the punctured formal neighborhood $\Spec(\widehat{K}_P)$
 allows us to study the ramification in a strictly local context.
 This restriction yields a $G$-torsor over a field, which necessarily decomposes into a disjoint union of spectra of finite separable field extensions
 $$ \cP|_{\Spec(\widehat{K}_P)} \cong \bigsqcup \Spec(F) $$
 such that:
 \begin{enumerate}
    \item These extensions $F/\widehat{K}_P$ are all isomorphic and Galois.
    \item The Galois group $H = \Gal(F/\widehat{K}_P)$ identifies with the stabilizer of any connected component of the pullback $\cP|_{\Spec(\widehat{K}_P)}$,
     which is a subgroup of $G$.
    \item  The number of such components is then given by the index $[G:H]$.
 \end{enumerate}

 \begin{definition}\label{definition:local_ramification_boundness}
    We say that the extension $F/\widehat{K}_P$ \textit{has ramification bounded by $r$} if the ramification group $H^r$ is trivial.
    We say \textit{$\cP$ has ramification bounded by $r$ at $P$} if the corresponding local field extensions $F/\widehat{K}_P$
    satisfy this condition.
 \end{definition}


\begin{definition}
    Let $\mdls = \sum n_i P_i$ be a modulus on $C$. We say the $G$-torsor $\cP \to U$ has
    \textbf{ramification bounded by $\mdls$} if, for each point $P_i$, the local restriction
    $\cP|_{\Spec(\widehat{K}_{P_i})}$ has ramification bounded by the coefficient $n_i$.
\end{definition}

This local property remains invariant under base change to the separable closure.
Let $\bar{s} = \Spec(\bar{k}) \to \Spec(k)$
be a geometric point.
The higher ramification groups for $\cP|_{\Spec(\widehat{K}_P)}$ are isomorphic to those of
the geometric restriction $\cP_{\bar{k}}$ at any point $\bar{x}$ lying over $P$.
Thus an equivalent definition:
\begin{definition}
    Let $\mdls = \sum n_i P_i$ be a modulus on $C$.
    The $G$-torsor $\cP \to U$ has \textbf{ramification bounded by $\mdls$}
    if for every geometric point $\bar{x}$ of $\mdls$ (with image $\bar{s}$ in $\Spec(k)$),
    the restriction of $\cP$ to$$ \Spec(\widehat{\mathcal{O}}_{C_{\bar{k}},\bar{x}}) \times_{C_{\bar{k}}} U_{\bar{k}} $$
    has ramification bounded by the multiplicity of $\mdls_{\bar{s}}$ at $\bar{x}$ in the sense of
    \Cref{definition:local_ramification_boundness}.
\end{definition}
These definitions are equivalent because $\widehat{\mathcal{O}}_{C_{\bar{k}},\bar{x}} \cong \widehat{\mathcal{O}}_{C,P} \otimes_k \bar{k}$. Furthermore, the notions of tame ramification and unramifiedness derived here coincide with the codimension-$1$ criteria in \Cref{definition:tamely_ramified_in_codim_1}.

There is another equivalent formulation through the language of characters.
The group $G$, being a finite abelian group over $k$, is \'etale.
Consequently, the correspondence between $G$-torsors and the fundamental group allows us to describe $G$-torsors
 via characters. Specifically, there is an isomorphism (\Cref{main:cor:abelian_equivilance_fundamental_torsors}):
\[ \mathrm{Hom}_{\mathrm{cont.}}(\pi_1^{\text{\'et}}(U, \overline{x}), G) \xrightarrow{\cong} \mathrm{Tors}(U_{\text{\'et}}, G) \]

In our local setting, where we consider the restriction to the complete local field $\widehat{K}_P$,
the fundamental group identifies with the absolute Galois group
$G_{\widehat{K}_P} := \mathrm{Gal}({\widehat{K}_P}^{\mathrm{sep}}/{\widehat{K}_P})$.
Thus, the restricted $G$-torsor $\mathcal{P}|_{\mathrm{Spec}(\widehat{K}_P)}$ corresponds to a unique continuous homomorphism:
\[ \rho : G_{\widehat{K}_P} \to G \]

Under this correspondence, the torsor $\mathcal{P}$ has ramification bounded by $r$ at $P$
if and only if the homomorphism $\rho$ trivializes the $r$-th ramification group in the upper numbering:
\[ \rho(G_{\widehat{K}_P}^r) = \{ 1 \} \]

\subsection*{Basic Properties of Ramification of \texorpdfstring{$G$-Torsors}{G-Torsors}}
We now prove some elementary lemmas regarding the ramification of $G$-torsors.
\begin{lemma}\label{lemma:bounding_ramification_of_contracted_product}
    Let $G$ be a finite abelian group and $X$ be a locally Noetherian scheme over a field $k$.
    Let $U \subset X$ be a dense open subset and let $(Z, \xi)$ be a prime divisor in the complement $X \setminus U$.
    Assume $\cP_1$ and $\cP_2$ are two $G$-torsors on $U_{\text{ét}}$, such that $\cP_1$ has ramification bounded by $r_1$ at $(Z, \xi)$, and $\cP_2$ has ramification bounded by $r_2$ at $(Z, \xi)$.
    Then their product $\cP_1 \otimes^G \cP_2$ has ramification bounded by $\max(r_1, r_2)$ at $(Z, \xi)$.
\end{lemma}
\begin{proof}
    Let $A=\widehat{\Oo_{X, \xi}}$ and let $K=\Frac(A)$.
    Let $\rho_1, \rho_2: G_K \to G$ be the associated continuous homomorphisms corresponding to the $G$-torsors
    $\cP_1|_{\Spec(K)}$, $\cP_2|_{\Spec(K)}$.
    We finish by noticing that the associated character of $(\cP_1 \otimes^G \cP_2)|_{\Spec(K)}$ is $\rho=\rho_1 + \rho_2$.
\end{proof}

\begin{lemma}\label{lemma:descend_ramification_along_etale}

Let $f: X \to Y$ be a morphism of locally Noetherian schemes. Let $U_X \subset X$ and $U_Y \subset Y$ be dense open subschemes such that $f^{-1}(U_Y) \subset U_X$.

Let $(Z_X, \xi_X)$ and $(Z_Y, \xi_Y)$ be prime divisors of $X$ and $Y$ such that $Z_X \cap U_X = \emptyset$ and $Z_Y \cap U_Y = \emptyset$.
Suppose $f(\xi_X) = \xi_Y$ and that $f$ is étale at $\xi_X$.
Let $\mathcal{P}$ be a $G$-torsor over $U_Y$, and let $f^{-1} \mathcal{P}$ be its pullback to $f^{-1}(U_Y) \subset U_X$.
Then, the ramification of $f^{-1} \mathcal{P}$ at $\xi_X$ is bounded by $r$ if and only if the ramification of $\mathcal{P}$ at $\xi_Y$ is bounded by $r$.
\end{lemma}
\begin{proof}
The boundedness of ramification is determined by the behavior of the torsor over the completion of the local rings at the generic points.
Let $A = \mathcal{O}_{Y, \xi}$ and $B = \mathcal{O}_{X, \xi_X}$ be the discrete valuation rings at the generic points, with fraction fields
$K$ and $L$ respectively. Since $f$ is étale at $\xi_X$, the map $A \to B$ is a flat, unramified (hence weakly unramified) local homomorphism.
Consequently, the extension of completions $\widehat{L} / \widehat{K}$ is a finite unramified extension of complete discretely valued fields.
We finish by noticing that the upper numbering filtration on the absolute Galois group is compatible with unramified base change.\footnote{
 Specifically, let $G_K = \operatorname{Gal}(K^{sep}/K)$ and $G_L = \operatorname{Gal}(L^{sep}/L)$. For an unramified extension, the Herbrand function is the identity, which implies that for any $r \geq 0$:
\[ G_L^r = G_K^r \cap G_L \]}
\end{proof}

% --- NEW (not in the current chapter): two further stability lemmas, needed
%     for the assembly of the general abelian case in 2.5.4. Their proofs are
%     character computations in the style of the two lemmas above. ---

\begin{lemma}\label{lemma:ramification_of_fiber_product}
    Let $G_1, G_2$ be finite abelian groups and let $X, U, (Z, \xi)$ be as in
    \Cref{lemma:bounding_ramification_of_contracted_product}. Let $\cP_1$ be a
    $G_1$-torsor and $\cP_2$ a $G_2$-torsor on $U_{\text{ét}}$, and consider the fiber
    product $\cP_1 \times_U \cP_2$, which is a $(G_1 \times G_2)$-torsor on
    $U_{\text{ét}}$. If $\cP_i$ has ramification bounded by $r_i$ at $(Z, \xi)$ for
    $i = 1, 2$, then $\cP_1 \times_U \cP_2$ has ramification bounded by
    $\max(r_1, r_2)$ at $(Z, \xi)$. Moreover, $\cP_1 \times_U \cP_2$ is tamely
    ramified (resp.\ unramified) at $(Z, \xi)$ if and only if both $\cP_1$ and
    $\cP_2$ are.
\end{lemma}
\begin{proof}
    Let $A=\widehat{\Oo_{X, \xi}}$, $K=\Frac(A)$, and let
    $\rho_i: G_K \to G_i$ be the characters of $\cP_i|_{\Spec(K)}$. The
    character of $(\cP_1 \times_U \cP_2)|_{\Spec(K)}$ is
    $\rho = (\rho_1, \rho_2) : G_K \to G_1 \times G_2$, since torsors under a
    product group correspond to pairs of characters,
    $H^1(K, G_1 \times G_2) = H^1(K, G_1) \times H^1(K, G_2)$, compatibly with
    the projections. Now for any closed subgroup $N \subset G_K$ (we take the
    upper-numbering group $G_K^r$, resp.\ the wild inertia subgroup, resp.\ the
    inertia subgroup) we have $\rho(N) = \{1\}$ if and only if
    $\rho_1(N) = \{1\}$ and $\rho_2(N) = \{1\}$.
\end{proof}

\begin{lemma}\label{lemma:ramification_of_induced_torsor}
    Let $\iota : H \hookrightarrow G$ be an injective homomorphism of finite abelian
    groups, let $X, U, (Z, \xi)$ be as above, let $\cQ$ be an $H$-torsor on
    $U_{\text{ét}}$, and let $\iota_* \cQ := (\cQ \times G)/H$ be the induced
    $G$-torsor (where $H$ acts by $h \cdot (q, g) = (qh,\ \iota(h)^{-1}g)$).
    Then $\iota_* \cQ$ has ramification bounded by $r$ at $(Z, \xi)$ if and only if
    $\cQ$ does; likewise for tame ramification and for unramifiedness.
\end{lemma}
\begin{proof}
    Let $A=\widehat{\Oo_{X, \xi}}$, $K=\Frac(A)$, and let $\chi: G_K \to H$ be the
    character of $\cQ|_{\Spec(K)}$; the character of $(\iota_*\cQ)|_{\Spec(K)}$ is
    $\iota \circ \chi$. Since $\iota$ is injective, $\iota \circ \chi$ and $\chi$
    have the same kernel, so for any closed subgroup $N \subset G_K$ we have
    $(\iota \circ \chi)(N) = \{1\}$ if and only if $\chi(N) = \{1\}$.
    (Concretely: the connected components of the two torsors over $\Spec(K)$ are
    spectra of the \emph{same} field extension $F/K$, the fixed field of
    $\ker \chi$; the induced torsor merely has $[G : H]$ times as many components.)
\end{proof}

% ---------------------------------------------------------------------------
%  2.3  Kummer, Artin--Schreier, and Artin--Schreier--Witt Theories
%  [Basic theory only, presented in parallel. Sources: the "Kummer and
%   Artin-Schreier Theories" block of LocalRamificationPreliminaries.tex;
%   the reduced-representative arguments extracted from the proofs of
%   lemma:realization_over_Gm (RamificationAfterBlowupIsTame.tex) and
%   prop:reduced_polynomial_representatives (RamificationAfterBlowupPPower.tex);
%   witt/WittVectorPrerequisites.tex is \input unchanged.]
% ---------------------------------------------------------------------------

\section{Kummer, Artin--Schreier, and Artin--Schreier--Witt Theories}\label{section:cyclic_extension_theories}

For cyclic extensions, Kummer, Artin--Schreier, and Artin--Schreier--Witt
theories provide explicit ways to classify extensions and calculate their
ramification. The three theories follow one and the same template. In each
case there is a short exact sequence of commutative group schemes over the
base field $K$,
\[
    1 \to \mu_e \to \bG_m \xrightarrow{(-)^e} \bG_m \to 1,
    \qquad
    0 \to \Z/p\Z \to \bG_a \xrightarrow{\wp} \bG_a \to 0,
    \qquad
    0 \to \Z/p^r\Z \to W_r \xrightarrow{\wp} W_r \to 0,
\]
whose long exact sequence in (\'etale or Galois) cohomology classifies the
$G$-torsors over $K$ --- for $G = \Z/e\Z$ (with $\mu_e \cong \Z/e\Z$ after a
choice of primitive root of unity, requiring $\gcd(e, \operatorname{char} K) = 1$
and $\mu_e \subset K$), $G = \Z/p\Z$, and $G = \Z/p^r\Z$ respectively --- by
an explicit quotient group:
\[
    K^\times/(K^\times)^e,
    \qquad
    K/\wp(K),
    \qquad
    W_r(K)/\wp(W_r(K)).
\]
Concretely, each class is realized by adjoining a solution of one explicit
equation: $X^e = a$, resp. $\wp(X) = X^p - X = a$, resp.
$\wp(\vec X) = \vec f$ in Witt coordinates.
When $K$ is moreover a local field $k((x))$, each class admits a
\emph{reduced representative} --- a Laurent polynomial (or vector of Laurent
polynomials) in $x^{-1}$ normalized so that the ramification of the extension
can be read off directly from the valuation: the ramification index in the
Kummer case, and the (highest) upper-numbering break in the Artin--Schreier
and Artin--Schreier--Witt cases. The comparison table at the end of the
section summarizes the three columns of this dictionary; the computations of
\Cref{section:symmetric_power_computations} will traverse each column in
turn.

Throughout this section, $K$ is a discrete valuation field with a perfect
residue field $\kappa$ of characteristic $p > 0$; from
\Cref{lemma:kummer_reduced_representatives} on, we further specialize to
$K = k((x))$ with $k$ algebraically closed.

\subsection{Kummer Theory}\label{subsection:kummer_theory}

Let $e$ be a positive integer prime to $p$, and assume $K$ contains the group
$\mu_e$ of $e$-th roots of unity (of order $e$, since $\gcd(e,p)=1$ and
$\kappa$ is perfect). Kummer theory states that every cyclic extension of
$K$ of degree dividing $e$ is generated by an $e$-th root: it is of the form
$K(a^{1/e})$ for some $a \in K^\times$, and the assignment
$a \mapsto K(a^{1/e})$ induces a bijection between the cyclic subgroups of
$K^\times/(K^\times)^e$ and the cyclic extensions of $K$ of degree dividing
$e$, with $\Gal(K(a^{1/e})/K)$ cyclic of order equal to the order of $a$ in
$K^\times/(K^\times)^e$, acting by multiplication by roots of unity on
$a^{1/e}$. The ramification of these extensions is computed by the valuation
of $a$:

\begin{theorem}[Ramification in Kummer Extensions, \cite{koch1997algebraic}, Proposition 1.83]\label{theorem:kummer_ramification}
Assume $K$ contains the $n$-th roots of unity $\mu_n$, with $\gcd(n, p) = 1$.
Let $L/K$ be the extension given by the equation $X^n = a$ for some
$a \in K^\times$ and denote by $G$ its Galois group. Then we have:

\begin{enumerate}
    \item If $v_K(a) \in n\mathbb{Z}$ and the image of $a\pi^{-v_K(a)}$ in the residue field $\kappa$ is an $n$-th power, the extension $L/K$ is trivial.
    \item If $v_K(a) \in n\mathbb{Z}$ and the image of $a\pi^{-v_K(a)}$ in the residue field $\kappa$ is not an $n$-th power, the extension $L/K$ is cyclic and unramified.
    \item If $v_K(a) \notin n\mathbb{Z}$, the extension $L/K$ is cyclic and ramified. Specifically, if $\gcd(|v_K(a)|, n) = 1$, the extension is totally ramified of degree $n$.
     Otherwise it has ramification index $\frac{n}{\gcd(|v_K(a)|, n)}$.
\end{enumerate}
Conversely, Kummer theory ensures that every cyclic extension of degree $n$, prime to $p$, of a field that contains the $n$-th roots of unity, is of the above form.
Moreover, in the above we can always take $a \in \Oo_K$.
\end{theorem}
Note that in the case of total ramification the extension is tamely ramified;
indeed, since $\gcd(n,p)=1$, \emph{every} extension appearing in
\Cref{theorem:kummer_ramification} is tamely ramified.

Over the local field $k((x))$ with $k$ algebraically closed, the classes
admit a particularly simple reduced representative --- a power of the
uniformizer:

\begin{lemma}[reduced representatives, Kummer case]\label{lemma:kummer_reduced_representatives}
    Let $k$ be algebraically closed of characteristic $p > 0$, let $K = k((x))$,
    and let $e$ be a positive integer prime to $p$. Then:
    \begin{enumerate}
        \item Every tamely ramified extension of $K$ is of the form
        $k((x^{1/e'}))$ for some $e'$ prime to $p$; consequently the tame
        quotient $G^t$ of $\Gal(K^{sep}/K)$ is procyclic,
        $G^t \cong \varprojlim_{p \nmid e'} \mu_{e'}(k) \cong \prod_{\ell \neq p} \Z_\ell$.
        \item Every class in $K^\times/(K^\times)^e$ is represented by $x^a$
        for a unique $a \in \{0, 1, \dots, e-1\}$; the classes of the $x^a$
        exhaust $\mathrm{Hom}_{\mathrm{cont.}}(G^t, \Z/e\Z) \cong \Z/e\Z$, and the class of
        $x^a$ is a surjective character (i.e., the extension $X^e = x^a$ is a
        field, totally ramified of degree $e$) precisely when $\gcd(a,e)=1$.
    \end{enumerate}
\end{lemma}
\begin{proof}
    (1) Since $k$ is algebraically closed, the residue field of $K$ is
    algebraically closed, so $K$ admits no non-trivial unramified extensions,
    and its tamely ramified extensions are exactly the Kummer extensions
    $k((x^{1/e'}))$ with $p \nmid e'$ (\Cref{theorem:kummer_ramification}).
    The description of $G^t$ follows by passing to the limit.

    (2) Since $k^\times$ is divisible and every unit of $k[[x]]$ with
    constant term $1$ is an $e$-th power (Hensel; $\gcd(e,p)=1$), we have
    $K^\times = x^{\Z} \cdot (K^\times)^e \cdot k^\times = x^{\Z}(K^\times)^e$,
    so the classes of $x^0, \dots, x^{e-1}$ exhaust $K^\times/(K^\times)^e
    \cong \Z/e\Z$. Under the Kummer isomorphism
    $K^\times/(K^\times)^e \cong \mathrm{Hom}_{\mathrm{cont.}}(G^t, \Z/e\Z)$ the class of
    $x^a$ is $a$ times a generator, hence surjective iff $\gcd(a,e)=1$; by
    \Cref{theorem:kummer_ramification}(3) this is exactly the totally
    ramified case of degree $e$.
\end{proof}

\subsection{Artin--Schreier Theory}\label{subsection:artin_schreier_theory}

For cyclic extensions of degree $p = \operatorname{char} \kappa$, the role of
the $e$-th power map is played by the \emph{Artin--Schreier operator}
$\wp(X) = X^p - X$, and the role of $K^\times/(K^\times)^e$ by the additive
quotient $K/\wp(K)$.

\begin{theorem}[Ramification in Artin--Schreier Extensions, \cite{Thomas2005}]\label{theorem:artin_schreier_ramification}
Let $\wp(x) = x^p - x$ be the Artin--Schreier operator.
Let $L/K$ be the extension given by the equation $X^p - X = a$ for some $a \in K$ and denote by $G$ its Galois group.
Then we have:
\begin{enumerate}
    \item If $v_K(a) > 0$ or if $v_K(a) = 0$ and $a \in \wp(K)$, the extension $L/K$ is trivial.
    \item If $v_K(a) = 0$ and if $a \notin \wp(K)$, the extension $L/K$ is cyclic of degree $p$ and unramified.
    \item If $v_K(a) = -m < 0$ with $m \in \mathbb{Z}_{>0}$ and if $m$ is prime to $p$, the extension $L/K$ is cyclic of degree $p$ again and totally ramified. Moreover,
     its ramification groups are given by:
     $$G = G^{(-1)} = \dots = G^{(m)} \quad \text{and} \quad G^{(m+1)} = 1.$$
\end{enumerate}
Conversely, Artin--Schreier theory ensures that every cyclic extension of degree $p$ takes this form.
\end{theorem}

The analogue of the representatives $x^a$ of
\Cref{lemma:kummer_reduced_representatives} is a polynomial in $x^{-1}$ whose
pole order is prime to $p$; the pole order is then exactly the unique
ramification break of \Cref{theorem:artin_schreier_ramification}(3).

\begin{lemma}[reduced representatives, Artin--Schreier case]\label{lemma:as_reduced_representatives}
    Let $k$ be algebraically closed of characteristic $p > 0$ and let
    $K = k((x))$. Every class in $K/\wp(K)$ admits a representative
    $f \in x^{-1}k[x^{-1}]$ which is either $0$ or of the form
    $$f = cx^{-m} + a_{-m+1}x^{-m+1} + \dots + a_{-1}x^{-1}, \qquad c \neq 0, \quad p \nmid m.$$
    In the latter case the extension $X^p - X = f$ is cyclic of degree $p$,
    totally ramified, with unique ramification break $m$.
\end{lemma}
\begin{proof}
    We show every class admits such a representative:
    \begin{itemize}
        \item $\wp$ is surjective on $k[[x]]$: given $h = \sum_{i \geq 0} b_i x^i$, solve
        $\wp(u) = h$ with $u = \sum_{i \geq 0} u_i x^i$ by choosing $u_0$ with $u_0^p - u_0 = b_0$ (possible as $k$ is algebraically closed)
        and setting recursively $u_i = u_{i/p}^p - b_i$, with the convention $u_{i/p} = 0$ when $p \nmid i$. Hence, writing an arbitrary
        element as $f^- + f^+$ with $f^- \in x^{-1}k[x^{-1}]$ and $f^+ \in k[[x]]$, we may
        choose a representative in $x^{-1}k[x^{-1}]$.
        \item If the representative has pole order $m$ divisible by $p$, with leading term $cx^{-m}$, then subtracting
        $\wp(c^{1/p}x^{-m/p}) = cx^{-m} - c^{1/p}x^{-m/p}$ strictly decreases the pole order while staying inside
        $x^{-1}k[x^{-1}]$. Iterating, we reach a representative $f$ which is either $0$ or has pole order $m$ prime to $p$.
    \end{itemize}
    The last assertion is \Cref{theorem:artin_schreier_ramification}(3),
    noting that a non-zero reduced $f$ has $v_K(f) = -m < 0$ with $p \nmid m$.
\end{proof}

Cyclic extensions of degree $p^r$ are classified by the same picture with the
additive group $\bG_a$ replaced by the ring scheme $W_r$ of Witt vectors of
length $r$, and $\wp = F - \mathrm{id}$ the Artin--Schreier--Witt operator.
The following subsection collects the required background on Witt vectors and
the Artin--Schreier--Witt classification
(\Cref{theorem:asw_theory}) together with the Schmid--Witt formula
(\Cref{theorem:schmid_witt_break_formula}) computing the ramification breaks
--- the degree-$p^r$ row of the dictionary.

% --- unchanged thesis file, read in place (defines subsection 2.3.3): ---
\providecommand{\bW}{\mathbb{W}}

\subsection{Artin--Schreier--Witt Theory}
\phantomsection\label{subsection:witt_vector_prerequisites}

Artin--Schreier theory, together with the ramification formula of
\Cref{theorem:artin_schreier_ramification}, gives complete control over the cyclic
extensions of degree $p$ of a local field of characteristic $p$. To treat the groups
$G = \Z/p^r\Z$ we need the analogous control over cyclic extensions of degree $p^r$,
which is provided by Witt's generalization: the field $K$ is replaced by the ring
$W_r(K)$ of $p$-typical Witt vectors of length $r$, the operator $\wp(u) = u^p - u$ by
its Witt-vector analogue $\wp = F - \mathrm{id}$, and the classical ramification-break
formula by the Schmid--Witt formula. The underlying Witt-ring algebra is collected in
\Cref{appendix:witt_vector_algebra}; this subsection retains the ASW classification,
the identities used in the proof, and the ramification statements. Our references for
the general theory are \cite[II, \S6]{serreLF} and \cite{Rabinoff2014Witt}.

Throughout, $p$ is a fixed prime, and $r \geq 1$ is an integer. All the Witt vectors
appearing in this thesis are $p$-typical and of finite length.


\paragraph{Working Witt-vector notation.}
For an \(\F_p\)-algebra \(A\), write \(W_r(A)=A^r\) for the ring of
\(p\)-typical Witt vectors of length \(r\), with Witt addition and subtraction
denoted by \(\oplus\) and \(\ominus\). We use the componentwise Frobenius \(F\),
the Verschiebung \(V\), and the truncation \(t\). The construction from Witt
polynomials, the ghost-map arguments, and the proofs of the identities among
these operators are given in \Cref{appendix:witt_vector_algebra}. In particular,
we use the triangular form of Witt addition
(\Cref{eq:witt_addition_shape}), the identities
\[
tW_r(A)=W_{r-1}(A),\qquad
\ker(t)=V^{r-1}A,\qquad
FV=VF=p,
\]
and the canonical identification \(W_r(\F_p)\cong\Z/p^r\Z\)
(\Cref{prop:witt_multiplication_by_p,example:witt_of_Fp}).

\begin{definition}\label{definition:asw_operator}
The \textbf{Artin--Schreier--Witt operator} on $W_r(A)$, $A$ an $\F_p$-algebra, is
\[
\wp \;:=\; F - \mathrm{id},
\qquad
\wp\,\vec u \;=\; F\vec u \ominus \vec u .
\]
For $r = 1$ this is the classical operator $\wp(u) = u^p - u$ of
\Cref{theorem:artin_schreier_ramification}.
\end{definition}

Since $F$ is a ring endomorphism, $\wp$ is an endomorphism of the additive group
$(W_r(A), \oplus)$:
$\wp(\vec u \oplus \vec v) = \wp\vec u \oplus \wp\vec v$. Moreover $\wp$ commutes with
$W_r(\sigma)$ for every ring homomorphism $\sigma$ (both $F$ and the ring operations
are natural), and with $V$ and $t$ by \Cref{prop:witt_multiplication_by_p}.

The following bookkeeping lemma is elementary but does the essential work of reducing
statements about $W_r$ to statements about $W_{r-1}$ and about $A$ itself: the
truncation $t$ and the inclusion $V^{r-1}$ of its kernel slice $W_r$ into lower-length
pieces, compatibly with $\wp$, and \emph{addition of an element of the kernel is
carry-free}.

\begin{lemma}[Witt bookkeeping]\label{lemma:witt_bookkeeping}
Let $A$ be an $\F_p$-algebra.
\begin{enumerate}[label=(\roman*)]
    \item $t \colon W_r(A) \to W_{r-1}(A)$ is a ring homomorphism commuting with $F$,
    hence with $\wp$; its kernel is $V^{r-1} A$.
    \item $V^{r-1} \colon A \to W_r(A)$ is additive, and
    $\wp \circ V^{r-1} = V^{r-1} \circ \wp$, where on the left $\wp$ is the operator
    on $W_r(A)$ and on the right the classical $\wp(c) = c^p - c$ on $A$.
    \item (carry-free top-level addition) For all $\vec a \in W_r(A)$ and $c \in A$,
    \[
    \vec a \oplus V^{r-1}(c) \;=\; (a_0, \dots, a_{r-2},\, a_{r-1} + c).
    \]
\end{enumerate}
\end{lemma}

\begin{proof}
(i) and (ii) restate parts of \Cref{prop:witt_multiplication_by_p} (for (ii): $V^{r-1}$
is additive and commutes with the componentwise $F$, hence with
$\wp = F - \mathrm{id}$).
(iii) By \Cref{lemma:ghost_principle} it suffices to verify the identity over
$\Z[a_0, \dots, a_{r-1}, c]$, where the ghost map is injective. Put
$\vec b := (0, \dots, 0, c)$ and $\vec s := \vec a \oplus \vec b$. For $n < r-1$ we
have $w_n(\vec b) = 0$, so $w_n(\vec s) = w_n(\vec a)$, and by induction on $n$
(solving for $s_n$, using that $p$ is a non-zero-divisor) $s_n = a_n$ for
$n \leq r-2$. For $n = r-1$, $w_{r-1}(\vec b) = p^{r-1} c$ gives
$p^{r-1} s_{r-1} = p^{r-1} a_{r-1} + p^{r-1} c$, i.e.\ $s_{r-1} = a_{r-1} + c$.
\end{proof}

\begin{rem*}
\Cref{lemma:witt_bookkeeping}(iii) fails for the addition of vectors supported in
lower positions: in general the carry polynomials $C_n$ of
\Cref{eq:witt_addition_shape} are non-zero, and Witt addition is \emph{not}
componentwise. In particular, operations performed componentwise on a Witt vector ---
such as discarding the regular parts of Laurent-series components --- do \emph{not}
in general respect Witt addition, and will be carefully avoided.
\end{rem*}

\subsubsection*{The Artin--Schreier--Witt isogeny and its torsors}

The componentwise ring operations make $\A^r_{\F_p}$ a ring scheme, the
\textbf{Witt vector group scheme} $\bW_r$: its underlying scheme is
$\A^r_{\F_p} = \Spec \F_p[T_0, \dots, T_{r-1}]$, with the group law given by the
addition polynomials $S_n$, so that $\bW_r(A) = W_r(A)$ functorially. The operator
$\wp = F - \mathrm{id}$ is an endomorphism of the group scheme $\bW_r$, and it is the
Witt-vector instance of a \emph{Lang isogeny}: by \Cref{example:witt_of_Fp} and the
lemma below, it sits in a short exact sequence of group schemes over $\F_p$ for the
\'etale topology
\begin{equation}\label{eq:asw_lang_sequence}
0 \longrightarrow \Z/p^r\Z \cong W_r(\F_p)
\longrightarrow \bW_r \xrightarrow{\;\wp\;} \bW_r \longrightarrow 0 .
\end{equation}

The corresponding covers admit the following completely explicit description, which
is the form in which we use them. For $\vec g \in W_r(A)$ we write
$A[\vec X]/(\wp \vec X \ominus \vec g)$, $\vec X = (X_0, \dots, X_{r-1})$, for the
quotient of $A[X_0, \dots, X_{r-1}]$ by the ideal generated by the $r$ components of
the Witt vector $\wp \vec X \ominus \vec g \in W_r(A[X_0, \dots, X_{r-1}])$.

\begin{lemma}[ASW covers are \'etale torsors]\label{lemma:asw_etale_torsor}
Let $A$ be an $\F_p$-algebra and $\vec g \in W_r(A)$. Let
\[
B \;:=\; A[\vec X]\big/\bigl(\wp \vec X \ominus \vec g\bigr).
\]
Then $B$ is finite free over $A$ with basis
$\bigl\{\prod_{n} X_n^{a_n} : 0 \le a_n \le p-1\bigr\}$, and \'etale over $A$. The
translation action of $W_r(\F_p) \cong \Z/p^r\Z$,
\[
m \colon \vec X \longmapsto \vec X \oplus \underline m ,
\]
makes $\Spec B$ a $\Z/p^r\Z$-torsor over $\Spec A$.
\end{lemma}

\begin{proof}
By \Cref{eq:witt_addition_shape} and the componentwise formula for $F$, the $n$-th
defining equation has the triangular form
\[
X_n^p - X_n - g_n + C_n'(X_0, \dots, X_{n-1};\, g_0, \dots, g_{n-1}) \;=\; 0
\]
for universal integral polynomials $C_n'$. Triangularity gives the free basis; the
Jacobian matrix of the system is lower triangular with diagonal entries
$p X_n^{p-1} - 1 = -1$, a unit, so $B$ is \'etale over $A$. Finally, $\Spec B$ is the
pullback along $\vec g \colon \Spec A \to \bW_r$ of the covering
$\wp \colon \bW_r \to \bW_r$; by the triangular form just established (over
$A = \F_p[T_0,\dots,T_{r-1}]$, $\vec g = (T_0,\dots,T_{r-1})$), $\wp$ is a surjective
\'etale homomorphism of group schemes, and its kernel is
$\{\vec u : F \vec u = \vec u\} = W_r(\F_p)$ since $F$ is componentwise. A surjective
\'etale group homomorphism is a torsor under its kernel, and torsors pull back to
torsors; hence $\Spec B$ is a $W_r(\F_p)$-torsor over $\Spec A$.
\end{proof}

\subsubsection*{Artin--Schreier--Witt theory for fields}

Over a field, the sequence \Cref{eq:asw_lang_sequence} computes all cyclic extensions
of $p$-power degree, exactly as classical Artin--Schreier theory does for degree $p$.

\begin{theorem}[Artin--Schreier--Witt theory; {\cite[\S3--4]{Thomas2005}}]
\label{theorem:asw_theory}
Let $M \supseteq \F_p$ be a field, $M^{s}$ a separable closure, and
$\Gamma = \Gal(M^{s}/M)$. Then $\wp$ is surjective on $W_r(M^{s})$ with kernel
$W_r(\F_p)$, and there is a canonical isomorphism of abelian groups
\[
\delta \colon\; W_r(M)\big/\wp\, W_r(M) \;\xrightarrow{\;\cong\;}\;
\Hom_{\mathrm{cont}}\bigl(\Gamma,\, \Z/p^r\Z\bigr) \;=\; H^1\bigl(M, \Z/p^r\Z\bigr),
\qquad
[\vec g] \longmapsto \bigl(\sigma \mapsto \sigma(\vec u) \ominus \vec u\bigr),
\]
where $\vec u \in W_r(M^{s})$ is any solution of $\wp \vec u = \vec g$. Under the
correspondence between characters and torsors
(\Cref{cor:abelian_equivilance_fundamental_torsors}), the class $[\vec g]$ corresponds
to the $\Z/p^r\Z$-torsor
\[
\Spec\, M[\vec X]\big/\bigl(\wp \vec X \ominus \vec g\bigr)
\]
of \Cref{lemma:asw_etale_torsor}. Moreover, for a class of order $p^s$ in
$W_r(M)/\wp W_r(M)$, the fixed field $M(\vec u) := M(u_0, \dots, u_{r-1})$ of
$\ker \delta[\vec g]$ is a cyclic extension of $M$ of degree $p^s$; the torsor is
connected if and only if $[\vec g]$ has order $p^r$, in which case
$\Spec M[\vec X]/(\wp \vec X \ominus \vec g) \cong \Spec M(\vec u)$ is (the spectrum
of) a cyclic field extension of degree $p^r$. Every cyclic extension of $M$ of degree
dividing $p^r$ arises this way.
\end{theorem}

\begin{proof}[Sketch of proof]
Surjectivity of $\wp$ on $W_r(M^{s})$ and the computation of its kernel follow from
the triangular equations in the proof of \Cref{lemma:asw_etale_torsor}: each component
of a preimage is obtained by solving a separable equation $X_n^p - X_n = (\dots)$, and
$F\vec u = \vec u$ holds if and only if $u_n^p = u_n$ for all $n$, i.e.\
$\vec u \in W_r(\F_p)$. Given $\vec g$ and a solution $\wp\vec u = \vec g$, for
$\sigma \in \Gamma$ we have
$\wp(\sigma \vec u \ominus \vec u) = \sigma \vec g \ominus \vec g = 0$ (using
\Cref{eq:witt_invariants_componentwise} and that $W_r(\sigma)$ commutes with $\wp$),
so $\sigma \vec u \ominus \vec u \in W_r(\F_p)$; one checks it is a continuous
homomorphism in $\sigma$ (the kernel is fixed pointwise by $\Gamma$), independent of
the choices modulo nothing and of the representative $\vec g$ modulo $\wp W_r(M)$.
Injectivity: if $\delta[\vec g] = 0$ then $\vec u$ is $\Gamma$-invariant, hence lies
in $W_r(M)$ by \Cref{eq:witt_invariants_componentwise}, so
$\vec g = \wp\vec u \in \wp W_r(M)$. Surjectivity of $\delta$ is the additive
Hilbert~90: $H^1(\Gamma, W_r(M^{s})) = 0$, by induction on $r$ along the exact
sequence $0 \to M^{s} \xrightarrow{V^{r-1}} W_r(M^{s}) \xrightarrow{t} W_{r-1}(M^{s})
\to 0$ of $\Gamma$-modules (\Cref{lemma:witt_bookkeeping}(i)) and the classical case
$H^1(\Gamma, M^{s}) = 0$. The remaining statements follow from
\Cref{theorem:equivalence_fundamental_covers} applied to the character
$\delta[\vec g]$, whose image has the same order as the class $[\vec g]$. For details
see \cite[\S3--4]{Thomas2005} or \cite{Rabinoff2014Witt}.
\end{proof}

The order of a class is partially visible in its $0$-th component:

\begin{cor}\label{cor:asw_class_order_and_first_component}
Let $M \supseteq \F_p$ be a field and $\vec g \in W_r(M)$ with $g_0 = 0$. Then
$p^{r-1}[\vec g] = 0$ in $W_r(M)/\wp W_r(M)$; equivalently, the cyclic extension
attached to $[\vec g]$ has degree at most $p^{r-1}$. Consequently, if $[\vec g]$ has
order $p^r$ (e.g.\ if $\vec g$ defines a connected torsor), then \emph{every}
representative of the class has non-zero $0$-th component.
\end{cor}

\begin{proof}
By \Cref{eq:witt_p_to_r_minus_1}, $p^{r-1} \vec g = (0, \dots, 0, g_0^{p^{r-1}}) = 0$
already in $W_r(M)$.
\end{proof}

\begin{rem*}\label{rem:asw_local_totally_ramified}
When $M = k((x))$ with $k$ algebraically closed, every finite separable extension
$L/M$ is totally ramified: the residue field extension is trivial (as $k$ is
algebraically closed) and $\Oo_L$ is a complete discrete valuation ring, so
$[L : M] = ef = e$. Thus a class of order $p^r$ in $W_r(k((x)))/\wp$ corresponds to a
cyclic, \emph{totally ramified} extension of degree $p^r$.
\end{rem*}



The isobaricity property of the Witt addition polynomials used in the later
degree estimate is recorded, with its proof, in
\Cref{lemma:witt_addition_isobaric}.

\subsubsection*{Ramification breaks: the Schmid--Witt formula}

Finally we state the generalization of the ramification formula of
\Cref{theorem:artin_schreier_ramification} to Artin--Schreier--Witt extensions. As in
the classical case, the formula reads off the largest upper-numbering ramification
break from the pole orders of a defining Witt vector --- provided the representative
is normalized so that no interference between the levels can occur. The relevant
normalization is the following.

\begin{definition}\label{definition:reduced_witt_vector}
Let $K = k((x))$. For $g \in K$ put $m(g) := \max\bigl(0, -v_x(g)\bigr)$ (the pole
order of $g$, or $0$ if $g$ is regular). A vector
$\vec g = (g_0, \dots, g_{r-1}) \in W_r(K)$ is \textbf{reduced} if for every $j$,
either $m(g_j) = 0$ or $p \nmid m(g_j)$.
\end{definition}

\begin{theorem}[Schmid--Witt break formula]\label{theorem:schmid_witt_break_formula}
Let $K = k((x))$ with $k$ perfect of characteristic $p$, and let $L/K$ be a cyclic
totally ramified extension of degree $p^r$ whose Artin--Schreier--Witt class
(\Cref{theorem:asw_theory}) is represented by a \emph{reduced} vector
$\vec g \in W_r(K)$, with $m_j := m(g_j)$. Then the largest upper-numbering
ramification break of $L/K$ equals
\[
u \;=\; \max_{0 \le j \le r-1} \; p^{\,r-1-j}\, m_j .
\]
\end{theorem}

This is due, for finite residue fields, to Kanesaka--Sekiguchi
\cite{KanesakaSekiguchi1979} and Thomas \cite{Thomas2005} (building on Schmid
\cite{Schmid1937} for $r = 1$), and for arbitrary perfect residue fields to
Elder--Keating \cite{ElderKeating2025}. We use it as stated, in one direction only:
\emph{evaluating the formula on one reduced representative computes $u$}; no
minimization over representatives is needed.

\begin{rem*}
Reducedness is exactly what makes the formula exact. If two levels $j < j'$ with
$m_j, m_{j'} \neq 0$ satisfied $p^{r-1-j} m_j = p^{r-1-j'} m_{j'}$, then
$m_{j'} = p^{\,j'-j} m_j$ would be divisible by $p$, contradicting
\Cref{definition:reduced_witt_vector}; hence the maximum is attained at a unique
level, and no cancellation of breaks between levels can occur.
\end{rem*}

\begin{rem*}
For $r = 1$, a reduced vector is a single element $g \in K$ whose pole order $m_0$ is
zero or prime to $p$, and \Cref{theorem:schmid_witt_break_formula} recovers
\Cref{theorem:artin_schreier_ramification}: the largest (indeed the only) upper break
is $m_0$. The existence of reduced representatives --- and, over $K = k((x))$, of
reduced representatives whose components are polynomials in $x^{-1}$ without constant
term --- is not part of the general theory quoted here; it will be proved where it is
used, by induction on $r$ through the truncation maps of
\Cref{lemma:witt_bookkeeping}.
\end{rem*}


\subsection{Reduced Representatives for Artin--Schreier--Witt Classes}\label{subsection:asw_reduced_representatives}

We now record the $W_r$-analogue of
\Cref{lemma:as_reduced_representatives}: every Artin--Schreier--Witt class
over $k((x))$ is represented by a vector of \emph{polynomials in $x^{-1}$}
with controlled pole orders. Recall from
\Cref{definition:reduced_witt_vector} the notion of a reduced vector; we add:

\begin{definition}\label{definition:reduced_polynomial_vector}
    A vector $\vec g = (g_0, \dots, g_{r-1}) \in W_r(k((x)))$ is a \textbf{reduced
    polynomial vector} if it is reduced (\Cref{definition:reduced_witt_vector}) and
    moreover every $g_j \in x^{-1}k[x^{-1}]$; then $m(g_j) = \deg_{x^{-1}} g_j$, and
    $m(g_j) = 0$ if and only if $g_j = 0$.
\end{definition}

\begin{lemma}[surjectivity of $\wp$ on integral Witt vectors]\label{lemma:wp_surjective_on_integral_witt}
    $\wp \colon W_r(k[[x]]) \to W_r(k[[x]])$ is surjective.
\end{lemma}

\begin{proof}
    Induction on $r$. For $r=1$ this is the first bullet in the proof of
    \Cref{lemma:as_reduced_representatives}: given $h = \sum_{i \geq 0} b_i x^i$, solve
    $u^p - u = h$ coefficientwise, using that $k$ is algebraically closed for the
    constant term. For $r>1$, let $\vec h \in W_r(k[[x]])$. By induction choose
    $\vec u' \in W_{r-1}(k[[x]])$ with $\wp \vec u' = t(\vec h)$, and lift it to
    $\tilde u := (\vec u', 0) \in W_r(k[[x]])$. By
    \Cref{lemma:witt_bookkeeping}(i),
    $t\bigl(\wp \tilde u \ominus \vec h\bigr) = \wp \vec u' \ominus t(\vec h) = 0$, so
    $\wp \tilde u \ominus \vec h = V^{r-1}(g)$ for some $g \in k[[x]]$. Choose
    $w \in k[[x]]$ with $\wp w = -g$ (case $r=1$ again). Then by
    \Cref{lemma:witt_bookkeeping}(ii) and additivity of $V^{r-1}$,
    \[
    \wp\bigl(\tilde u \oplus V^{r-1}(w)\bigr)
    = \wp \tilde u \oplus V^{r-1}(\wp w)
    = \bigl(\vec h \oplus V^{r-1}(g)\bigr) \oplus V^{r-1}(-g) = \vec h. \qedhere
    \]
\end{proof}

\begin{prop}[reduced polynomial representatives]\label{prop:reduced_polynomial_representatives}
    Every class in $W_r(k((x)))/\wp W_r(k((x)))$ has a reduced polynomial
    representative: for every $\vec f \in W_r(k((x)))$ there exists
    $\vec u \in W_r(k((x)))$ such that $\vec g := \vec f \ominus \wp \vec u$ has all
    components in $x^{-1}k[x^{-1}]$ with pole orders either $0$ or prime to $p$.
\end{prop}

\begin{proof}
    Induction on $r$.

    \emph{Base case $r=1$.} This is \Cref{lemma:as_reduced_representatives}:
    first, writing $f = f^- + f^+$ with
    $f^- \in x^{-1}k[x^{-1}]$ and $f^+ \in k[[x]]$, surjectivity of $\wp$ on $k[[x]]$
    removes $f^+$; then, as long as the leading term is $cx^{-m}$ with $p \mid m$,
    $m>0$, subtracting $\wp(c^{1/p}x^{-m/p}) = cx^{-m} - c^{1/p}x^{-m/p}$ strictly
    decreases the pole order while staying inside $x^{-1}k[x^{-1}]$; this terminates
    at $0$ or at pole order prime to $p$.

    \emph{Inductive step.} Let $\vec f \in W_r(k((x)))$. By induction there is
    $\vec u' \in W_{r-1}(k((x)))$ such that
    $\vec g' := t(\vec f) \ominus \wp \vec u'$ is a reduced polynomial vector in
    $W_{r-1}$. Lift $\vec u'$ to $\tilde u := (\vec u', 0) \in W_r(k((x)))$ and set
    $\vec h := \vec f \ominus \wp \tilde u$. By \Cref{lemma:witt_bookkeeping}(i),
    $t(\vec h) = \vec g'$, i.e.\
    \[
    \vec h = (g'_0, \dots, g'_{r-2},\, h_{r-1}), \qquad h_{r-1} \in k((x)),
    \]
    with the first $r-1$ components already reduced polynomials. It remains to repair
    the last component \emph{without touching the others}; by
    \Cref{lemma:witt_bookkeeping}(iii), adjustments of the form
    $\ominus\, \wp\bigl(V^{r-1}(z)\bigr) = \ominus\, V^{r-1}(\wp z)$ change only the
    last component, by $-\wp(z)$. So:
    \begin{itemize}
        \item Write $h_{r-1} = h^- + h^+$ with $h^- \in x^{-1}k[x^{-1}]$,
        $h^+ \in k[[x]]$, choose $w \in k[[x]]$ with $\wp w = h^+$
        (\Cref{lemma:wp_surjective_on_integral_witt} with $r=1$), and replace
        $\vec h$ by $\vec h \ominus \wp(V^{r-1}w)$: the last component becomes $h^-$.
        \item While the last component has leading term $cx^{-m}$ with $p \mid m$,
        $m>0$, replace $\vec h$ by
        $\vec h \ominus \wp\bigl(V^{r-1}(c^{1/p}x^{-m/p})\bigr)$: the last component
        changes by $-\wp(c^{1/p}x^{-m/p}) = -cx^{-m} + c^{1/p}x^{-m/p}$, so its pole
        order strictly decreases and it stays in $x^{-1}k[x^{-1}]$. This terminates at
        $0$ or at pole order prime to $p$.
    \end{itemize}
    The result is a reduced polynomial vector $\wp$-equivalent to $\vec f$.
\end{proof}

\begin{rem*}
    The individual components of a Witt vector cannot be truncated to their principal
    parts one at a time --- Witt addition is not componentwise
    (\Cref{eq:witt_addition_shape}), so a componentwise operation changes the
    Artin--Schreier--Witt class. The point of
    \Cref{prop:reduced_polynomial_representatives} is that all the adjustments are
    performed by subtracting genuine elements of $\wp W_r(k((x)))$, inside the Witt
    ring.
\end{rem*}

\subsection*{The Three Theories Side by Side}

The following table summarizes the parallel structure of the three theories
over $K = k((x))$, $k$ algebraically closed of characteristic $p > 0$. Each
column supplies exactly the ingredients that the corresponding computation in
\Cref{section:symmetric_power_computations} consumes: the explicit equation
(to realize a local torsor over $\bG_m$), the reduced representative (to
control pole orders), and the ramification criterion (to read off the
conclusion at the generic point of the exceptional divisor).

\begin{center}
\renewcommand{\arraystretch}{1.6}
\begin{tabular}{|l|c|c|c|}
\hline
 & \textbf{Kummer} & \textbf{Artin--Schreier} & \textbf{Artin--Schreier--Witt} \\
\hline
group $G$ & $\Z/e\Z$, $\gcd(e,p)=1$ & $\Z/p\Z$ & $\Z/p^r\Z$ \\
\hline
ambient group scheme & $\bG_m$ & $\bG_a$ & $W_r$ \\
\hline
isogeny & $X \mapsto X^e$ & $\wp(X) = X^p - X$ & $\wp(\vec X) = F\vec X \ominus \vec X$ \\
\hline
defining equation & $X^e = a$ & $X^p - X = a$ & $\wp(\vec X) = \vec f$ \\
\hline
classifying group & $K^\times/(K^\times)^e$ & $K/\wp(K)$ & $W_r(K)/\wp W_r(K)$ \\
\hline
reduced representative & $x^a$, $0 \le a < e$ &
\makecell{$f \in x^{-1}k[x^{-1}]$, \\ $p \nmid m = $ pole order} &
\makecell{$\vec f = (f_0, \dots, f_{r-1})$, \\ $f_j \in x^{-1}k[x^{-1}]$, $p \nmid m_j$} \\
\hline
classification & \Cref{lemma:kummer_reduced_representatives} & \Cref{lemma:as_reduced_representatives} & \Cref{theorem:asw_theory}, \Cref{prop:reduced_polynomial_representatives} \\
\hline
ramification & \makecell{tame; ram.\ index \\ $e/\gcd(a, e)$} &
\makecell{totally ramified; \\ unique break $m$} &
\makecell{totally ramified (if $p \nmid m_0$, $f_0 \neq 0$); \\ largest break $\max_j p^{\,r-1-j} m_j$} \\
\hline
criterion & \Cref{theorem:kummer_ramification} & \Cref{theorem:artin_schreier_ramification} & \Cref{theorem:schmid_witt_break_formula} \\
\hline
unramified iff & $e \mid a$ & $f = 0$ & $\vec f = \vec 0$ \\
\hline
\end{tabular}
\end{center}

% ---------------------------------------------------------------------------
%  2.4  Prelude: shared geometric background, each item stated ONCE.
%  [Sources: the blowup-of-affine-space computation appeared three times in
%   RamificationAfterBlowupIsTame.tex (lines 629-649, 868-889 and inside
%   lemma:formal_local_invariance); the normal-model discussion and
%   lemma:torsor_unramified_codim_one from lines 124-204; the formal-local
%   invariance block from lines 399-561; lemma:symmetric_regularity from
%   RamificationAfterBlowupPPower.tex; the invariant-theory scaffold is
%   factored out of the proofs of theorem:ramification_after_blowup_P1 and
%   prop:witt_cover_identification.]
% ---------------------------------------------------------------------------

\section{Prelude: The Local Geometry of Blowups of Symmetric Powers}\label{section:blowup_prelude}

This section collects the geometric ingredients that the three computations
of \Cref{section:symmetric_power_computations} share: the complete local
field at the generic point of the exceptional divisor, the formal--local
invariance of the whole problem (which reduces us to $\bP^1$ and to explicit
torsors over $\bG_m$), the invariant-theoretic description of the
symmetric-power cover, and a regularity criterion for symmetric functions.
Each is stated once and reused verbatim in all three cases.

\subsection*{The Blowup of Affine Space at the Origin}

\begin{prop}\label{prop:blowup_local_ring_at_exceptional}
Let $X = \mathbb{A}_k^n$ be the affine $n$-space over a field $k$, let $0 \in X$ be the origin,
and let $\tilde{X} = \Bl_0(X)$ be the blowup of $X$ at the origin, so that
$\tilde{X} \subset X \times_k \mathbb{P}^{n-1}_k$ is defined by the equations $x_i u_j = x_j u_i$, where $[u_1 : \dots : u_n]$ are the homogeneous coordinates of
$\mathbb{P}^{n-1}_k$. Let $E \subset \tilde{X}$ be the exceptional divisor --- the fiber over the origin,
$E = \{ (0, [u_1 : \dots : u_n]) \}$, of codimension 1 in $\tilde{X}$ --- let
$\eta \in E$ be its generic point, and let $R = \mathcal{O}_{\tilde{X}, \eta}$. Then:
\begin{enumerate}
    \item $R$ is a discrete valuation ring with fraction field $K = K(X) = k(x_1, \dots, x_n)$.
    On the affine chart $U_1$ where $u_1 \neq 0$, we have $x_i = \frac{u_i}{u_1} x_1$, the coordinate ring is
    $\mathcal{O}_{\tilde{X}}(U_1) = k\left[x_1, \frac{u_2}{u_1}, \dots, \frac{u_n}{u_1}\right]$,
    and $\eta$ corresponds to the prime ideal $(x_1)$, so
    $R = k[x_1, \frac{u_2}{u_1}, \dots, \frac{u_n}{u_1}]_{(x_1)}$.
    \item The residue field is $\kappa(\eta) = k(\frac{u_2}{u_1}, \dots, \frac{u_n}{u_1})$, and the completion of $R$ is
    $\hat{R} = \kappa(\eta)[[x_1]]$. The element $x_1$ is a uniformizer; so is every $x_i$, since
    $x_i = (\frac{u_i}{u_1}) x_1$ with $\frac{u_i}{u_1}$ a unit in $R$.
    \item The valuation $\nu_E$ of $R$ computes the order of vanishing at the origin:
    for a monomial $M = x_1^{a_1} \dots x_n^{a_n}$ we have
    $M = x_1^{\sum a_i} \left( \frac{u_2}{u_1} \right)^{a_2} \dots \left( \frac{u_n}{u_1} \right)^{a_n}$
    with the parenthesized factor a unit in $R$, so $\nu_E(M) = \sum a_i = \deg M$; consequently, for any polynomial
    $f = f_d + f_{d+1} + \dots + f_l$ with $f_i$ homogeneous of degree $i$ and $f_d \neq 0$, we have $\nu_E(f) = d$.
\end{enumerate}
\end{prop}

The case relevant to us is the symmetric power: the situation is identical
after replacing the coordinates $x_1, \dots, x_d$ by the elementary symmetric
polynomials. Let $e_1, \dots, e_d$ be the elementary symmetric polynomials in
$x_1, \dots, x_d$, so that $C^{(d)} = \Sym^d_k(\mathbb{A}^1) = \Spec k[e_1, \dots, e_d] \cong \mathbb{A}^d_k$
locally at the origin. Blowing up the origin $Z = V(e_1, \dots, e_d)$ and applying
\Cref{prop:blowup_local_ring_at_exceptional} in the coordinates $e_1, \dots, e_d$ (on the chart $u_d \neq 0$), the local
ring at the generic point $\eta$ of the exceptional divisor is
\[
R = k\left[e_d, \tfrac{u_1}{u_d}, \dots, \tfrac{u_{d-1}}{u_d}\right]_{(e_d)},
\qquad e_i = \tfrac{u_i}{u_d}\, e_d \ \text{ all uniformizers},
\qquad \tfrac{e_i}{e_d} = \tfrac{u_i}{u_d} \in \kappa(\eta),
\]
with residue field $\kappa(\eta) = k(\frac{e_1}{e_d}, \dots, \frac{e_{d-1}}{e_d})$,
completion $\hat{R} = \kappa(\eta)[[e_d]]$, and complete local field

\begin{equation}\label{eq:complete_field_at_generic_point}
     \widehat{K}_\eta = k\left(\tfrac{e_1}{e_d}, \dots, \tfrac{e_{d-1}}{e_d}\right)((e_d)).
\end{equation}

\subsection*{Normal Models of Symmetric Powers of Torsors}

Let $C$ be a smooth, proper (projective) curve over a field $k$, and let $U \subset C$ be a dense open subset.
Let $G$ be a finite abelian group, and let $\pi_U: U' \to U$ be a $G$-torsor in $U_{\text{\'Et}}$.
Let $C'$ be the normal proper model of $U'$ over $k$ (i.e. the unique smooth proper curve which has $U'$ as a dense open subset).
Let $\pi: C' \to C$ be the induced morphism (which is finite, generically \'etale over $U$).
Let $U'^{(d)}$ be the standard $d$-th symmetric product of $U'$ as a scheme, and let $U'^{[d]}$ be the torsor-theoretic symmetric product;
by \Cref{main:cor:sym-power-quotient} we have a canonical morphism $q_U: U'^{(d)} \to U'^{[d]}$.
Note that $C'^{(d)}$ is smooth and finite over $C^{(d)}$, hence it is the normalization of $C^{(d)}$ inside $K(C'^{(d)})$.
$U'^{(d)}$ is a dense open subset of $C'^{(d)}$, hence $C'^{(d)}$ is a normal proper model of $U'^{(d)}$.
Let $C'^{[d]}$ be the normalization of $C^{(d)}$ in the function field $K(U'^{[d]})$;
then $q_U$ extends to a map $q: C'^{(d)} \to C'^{[d]}$, and we have the following diagrams:
\[
\begin{tikzcd}
	U' \arrow[r, hook] \arrow[d, "\pi_U"'] & C' \arrow[d, "\pi"] \\
	U \arrow[r, hook] & C
\end{tikzcd}
\qquad \qquad
\begin{tikzcd}
	{U'^{(d)}} \arrow[r, hook] \arrow[d, "q_U"'] & {C'^{(d)}} \arrow[d, "q"] \\
	{U'^{[d]}} \arrow[r, hook] \arrow[d] & {C'^{[d]}} \arrow[d] \\
	{U^{(d)}} \arrow[r, hook] & {C^{(d)}}
\end{tikzcd}
\]

\begin{lemma}\label{lemma:torsor_unramified_codim_one}
In the context of the preceding discussion, the morphism $U'^{(d)} \to U'^{[d]} \hookrightarrow C'^{[d]}$
is unramified in codimension $1$ (see \Cref{definition:tamely_ramified_in_codim_1}).
\end{lemma}
\begin{proof}

Let $C' \setminus U' = \{x_1, \ldots, x_m \}$. For each $x \in C' \setminus U'$, define the locus
$$Z_x = \{D \in C'^{(d)} \mid x \in \mathrm{Supp}(D) \}.$$
Each $Z_x$ is a prime divisor in the symmetric product $C'^{(d)}$, and the boundary $C'^{(d)} \setminus U'^{(d)}$ is precisely the finite union $\bigcup_{i=1}^m Z_{x_i}$.

Let $\eta$ be the generic point of $Z_x$.
It suffices to show that the local ring extension $\mathcal{O}_{C'^{[d]}, q(\eta)} \hookrightarrow \mathcal{O}_{C'^{(d)}, \eta}$ is weakly unramified (\Cref{definition:weakly_unramified_and_tamely_ramified_DVR}).

Let $p: C'^d \to C'^{(d)}$ be the natural quotient. Consider the closed subvariety
$$H = \{x\} \times C' \times \cdots \times C' \subset C'^d.$$

The subvariety $H$ is irreducible, and its image under $p$ is exactly $Z_x$.
Consequently, $p$ maps the generic point $\eta_H$ of $H$ to the generic point $\eta$ of $Z_x$.
\[
\begin{tikzcd}
\eta_H \arrow[d] \arrow[r] & H \arrow[r] \arrow[d] & C'^d \arrow[d, "p"] \\
\eta \arrow[r] \arrow[d] & Z_x \arrow[r] \arrow[d] & C'^{(d)} \arrow[d, "q"] \\
q(\eta) \arrow[r] & q(Z_x) \arrow[r] & C'^{[d]}
\end{tikzcd}
\]
It suffices to show that the local ring extension $\mathcal{O}_{C'^{[d]}, q(\eta)} \hookrightarrow \mathcal{O}_{C'^d, \eta_H}$ is weakly unramified.
Let $R = \mathcal{O}_{C'^d, \eta_H}$, and let $A=\mathcal{O}_{C'^{[d]}, q(\eta)}$.

Recall that $C'^{[d]}$ is the quotient of $C'^d$ by the finite group $\Xi = \KG \rtimes S_d$ (see \Cref{main:cor:sym-power-quotient}),
and the morphism $q \circ p : C'^d \to C'^{[d]}$ is a \textit{branched Galois cover} (finite, surjective map of normal integral schemes $X \to Y$, with function field extension $K(X)/K(Y)$ Galois).

Hence the decomposition group of $\eta_H$ in $\Xi$ and the inertia group of $\eta_H$ in $\Xi$ are defined:
$$D_{\eta_H} = \{ g \in \Xi \mid g(\eta_H) = \eta_H \}$$
$$I_{\eta_H} = \{ g \in D_{\eta_H} \mid g^*(u) \equiv u \pmod{\mathfrak{m}_{R}} \text{ for all } u \in R \}$$

\textbf{Claim: } The inertia group $I_{\eta_H}$ is trivial.

Let $K=K(C')$ be the function field of $C'$; then
$K(C'^d) = K^{\otimes_k d} \cong \Frac(K_1 \otimes_k K_2 \otimes_k \cdots \otimes_k K_d)$
where $K_i$ is the copy of $K$ corresponding to the $i$-th factor of $C'^d$.
The closed subscheme $H$ corresponds to the prime ideal $\pp_H \subset \Oo_{C'^d}$, and $R=\Oo_{C'^d, \pp_H}$.
The residue field of $R$ is exactly $K(H) \cong \Frac(k(x) \otimes_k K_2 \otimes_k \cdots \otimes_k K_d)$ where $k(x)$ is the residue field of $x$.

The action of $\gamma = (g_1, \ldots, g_d, \sigma) \in \Xi$, via the pullback $(\gamma^*)$ on $K(C'^d)$ does two things (in that order):
\begin{enumerate}
    \item It permutes the tensor factors according to $\sigma$. A function strictly in $K_i$ is mapped to a function in $K_{\sigma(i)}$.
    \item It applies the field automorphism $g_i^* \in \mathrm{Aut}(K/k)$ to the respective factors.
\end{enumerate}
Let $\gamma \in I_{\eta_H} \subset \Xi$. Then the induced automorphism $\bar{\gamma}^*$ on $K(H)$ must be the identity map.
This implies that $\sigma$ must fix $i$ for every $i \geq 2$ (the $K_i$ are linearly independent), and hence $\sigma=id$.
The trivial action of $\bar{\gamma}^*$ restricts to a trivial action of $g_i^*$ on $K_{i}$ for every $i \geq 2$, and hence
$g_i = 1$ for every $i \geq 2$ (since the action of $G$ on $K(C')$ is faithful).
Finally, since $g_1 g_2 \cdots g_d = 1$, we also have $g_1 = 1$.
Since the inertia group $I_{\eta_H}$ is trivial, the local ring extension $A \hookrightarrow R$ is weakly unramified.
(See for example \stackstag{09E3}, \stackstag{0BTF})

\end{proof}

\subsection*{Formal--Local Invariance and Reduction to \texorpdfstring{$\bP^1$}{P1}}

The next two results make precise the statement that \emph{the ramification of
$\cP^{[d]}$ at $\eta_\ndls$ is a formal--local invariant of $\cP$ at $P$}:
it depends only on the completed local ring $\widehat{\Oo_{C,P}} \cong k[[x]]$ together with the restriction of $\cP$
to the punctured formal disc $\Spec(k((x)))$, since all of the constructions involved --- self products, quotients by finite
groups, and blowups --- commute with the flat base change to the completion; for blowups, this is
\Cref{main:theorem:blowup_flat_base_change}.
We emphasize that this is not a general principle that can simply be quoted: it is a statement about the particular
constructions at hand, and it is proved by checking each construction separately.

Throughout the rest of this section we assume that $k$ is algebraically closed; in particular every closed point of $C$ is $k$-rational, and $\ndls = dP$ with $P \in C(k)$ and $d = \deg \ndls$.

\begin{remnum}\label{rem:algebraically_closed_harmless}
    The assumption that $k$ is algebraically closed is harmless in our situation, for the following reasons.
    \begin{enumerate}
        \item \textbf{(Base change to $\bar{k}$.)} Assume $k$ is perfect, so that $\bar{k}=k^{sep}$.
        The formation of $C^{(d)}$ (a quotient by a finite group), of the blowup $X_\ndls$ along $Z_\ndls$, and of
        $\cP^{[d]}$ (\Cref{main:cor:sym-power-quotient}) commutes with the flat base change $\Spec(\bar{k}) \to \Spec(k)$.
        Since $E_\ndls \cong \bP^{\deg \ndls - 1}_k$ is geometrically irreducible, $(E_\ndls)_{\bar{k}}$ is irreducible
        and its generic point $\bar{\eta}$ lies over $\eta_\ndls$; moreover
        $\Oo_{X_\ndls, \eta_\ndls} \to \Oo_{(X_\ndls)_{\bar{k}}, \bar{\eta}}$ is a weakly unramified extension of discrete
        valuation rings with separable residue field extension. The proof of \Cref{lemma:descend_ramification_along_etale}
        uses only these two properties of the extension (there the extension was moreover finite étale), so it applies
        verbatim: the ramification of $(\cP_{\bar{k}})^{[d]} = (\cP^{[d]})_{\bar{k}}$ at $\bar{\eta}$ is bounded by $r$ if
        and only if the ramification of $\cP^{[d]}$ at $\eta_\ndls$ is. Hence unramifiedness and tameness at $\eta_\ndls$
        may be checked after base change to $\bar{k}$.
        \item \textbf{(Rationality of $P$ --- a caveat.)} If $P$ is a closed point with $\kappa(P) \neq k$, then after base
        change to $\bar{k}$ the point $Z_{\ndls}$ corresponds to the divisor $\sum_{\bar{P} \mapsto P} n \bar{P}$ on
        $C_{\bar{k}}$, whose support consists of $[\kappa(P):k]$ distinct points, whereas the computation below treats a
        modulus supported at a single rational point. We therefore prove \Cref{theorem:torsors_after_blowup_is_tame} under
        the standing assumption that $k$ is algebraically closed; this suffices for our purposes, since in the proof of
        \Cref{main:theorem:SymmetricPowerOfSheafIsTamelyRamified} the field is assumed algebraically closed before the present
        theorem is invoked.
    \end{enumerate}
\end{remnum}

We now formulate the formal--local invariance precisely. Fix $P \in C(k)$ lying in the support of $\mdls$, set $\ndls = dP$,
and fix a uniformizer $x$ of $\Oo_{C,P}$; it induces an isomorphism $\widehat{\Oo_{C,P}} \cong k[[x]]$.
The scheme $\Spec(k[[x]])$ has two points --- the closed point, mapping to $P$, and the generic point, mapping to the generic
point of $C$ --- so the punctured formal disc $\Spec(k((x)))$ maps to $U$, and we denote by
$$\cP_x := \cP \times_U \Spec(k((x)))$$
the corresponding restriction of $\cP$; it is a $G$-torsor over $\Spec(k((x)))$.
Let $e_1, \dots, e_d \in k[[x_1, \dots, x_d]]$ be the elementary symmetric polynomials in $x_1, \dots, x_d$, and set
$$ V = \Spec\left(k[[e_1, \dots, e_d]][e_d^{-1}]\right), \qquad W = \Spec\left(k[[x_1, \dots, x_d]][(x_1 \cdots x_d)^{-1}]\right).$$
The inclusion $k[[e_1, \dots, e_d]] \subset k[[x_1, \dots, x_d]]$ induces a finite flat morphism $\hat{r}: W \to V$
(note that inverting $e_d = x_1 \cdots x_d$ inverts each $x_i$), and for each $i$ we let
$$ q_i : W \to \Spec(k((x))), \qquad x \mapsto x_i.$$
The group $S_d$ acts on $W$ over $V$ by permuting $x_1, \dots, x_d$, and $q_{\sigma(i)} = q_i \circ \sigma$.

\begin{lemma}[Formal--local invariance]\label{lemma:formal_local_invariance}
    In the above notation:
    \begin{enumerate}
        \item The choice of $x$ induces an isomorphism $\widehat{\Oo_{C^{(d)}, Z_\ndls}} \cong k[[e_1, \dots, e_d]]$ carrying
        the maximal ideal to $(e_1, \dots, e_d)$, such that the induced morphism $c: \Spec(k[[e_1, \dots, e_d]]) \to C^{(d)}$
        is flat and satisfies $c^{-1}(U^{(d)}) = V$; we write $c_V : V \to U^{(d)}$ for the restriction.
        \item There is an induced isomorphism $\widehat{\Oo_{X_\ndls, \eta_\ndls}} \cong \kappa[[e_d]]$, where
        $\kappa = k(\frac{e_1}{e_d}, \dots, \frac{e_{d-1}}{e_d})$. In particular
        $\widehat{K}_{\eta_\ndls} := \Frac(\widehat{\Oo_{X_\ndls, \eta_\ndls}}) \cong \kappa((e_d))$, and the canonical
        morphism $\Spec(\widehat{K}_{\eta_\ndls}) \to U^{(d)}$ factors as
        $\Spec(\kappa((e_d))) \to V \xrightarrow{c_V} U^{(d)}$, where the first map does not depend on $(C, \cP)$.
        \item There is a canonical isomorphism of $G$-torsors over $V$
        $$ \cP^{[d]} \times_{U^{(d)}} V \;\cong\; \Xi \backslash \left( q_1^{-1}\cP_x \times_W \cdots \times_W q_d^{-1}\cP_x \right).$$
    \end{enumerate}
    Consequently, the $G$-torsor $\cP^{[d]}|_{\Spec(\widehat{K}_{\eta_\ndls})}$ --- and hence whether the ramification of
    $\cP^{[d]}$ at $\eta_\ndls$ is bounded by $r$ (resp.\ tame, resp.\ unramified) --- depends only on $d$ and on the
    isomorphism class of the $G$-torsor $\cP_x$ over $k((x))$, and not otherwise on the pair $(C, \cP)$.
\end{lemma}

\begin{proof}
    \textbf{(1)} Since $P$ is a $k$-rational point of the smooth curve $C$, the uniformizer $x$ induces
    $\widehat{\Oo_{C,P}} \cong k[[x]]$. As completion at rational points takes products of $k$-varieties to completed tensor
    products, we get $\widehat{\Oo_{C^d, (P, \dots, P)}} \cong k[[x]] \widehat{\otimes}_k \cdots \widehat{\otimes}_k k[[x]] = k[[x_1, \dots, x_d]]$.
    Write $A = \Oo_{C^{(d)}, Z_\ndls}$ and $B = \Oo_{C^d, (P, \dots, P)}$.
    Since $k$ is algebraically closed, a closed point $(Q_1, \dots, Q_d) \in C^d$ maps to $Z_\ndls = dP$ under the finite
    projection $r: C^d \to C^{(d)}$ if and only if $Q_1 = \dots = Q_d = P$; hence $(P, \dots, P)$ is the unique point of
    $C^d$ over $Z_\ndls$, $B = (r_* \Oo_{C^d})_{Z_\ndls}$ is a finite $A$-module, and $A = B^{S_d}$ (formation of invariants
    commutes with localization at invariant primes).
    Since $\mathfrak{m}_B$ is the unique maximal ideal of $B$ over $\mathfrak{m}_A$, we have
    $\mathfrak{m}_B^N \subset \mathfrak{m}_A B \subset \mathfrak{m}_B$ for some $N$, so the $\mathfrak{m}_A$-adic and
    $\mathfrak{m}_B$-adic topologies on $B$ agree, and therefore $B \otimes_A \hat{A} \cong \hat{B} \cong k[[x_1, \dots, x_d]]$.
    Now $A = B^{S_d}$ is the kernel of $B \to \prod_{\sigma \in S_d} B$, $b \mapsto (\sigma(b) - b)_\sigma$; as $\hat{A}$ is
    flat over $A$, tensoring with $\hat{A}$ preserves this kernel, so
    $$\hat{A} \cong (B \otimes_A \hat{A})^{S_d} \cong k[[x_1, \dots, x_d]]^{S_d}.$$
    Finally, $k[[x_1, \dots, x_d]]^{S_d} = k[[e_1, \dots, e_d]]$: the ring $k[x_1, \dots, x_d]$ is a finite free module over
    $k[e_1, \dots, e_d]$, the same topological argument identifies $k[[x_1, \dots, x_d]]$ with
    $k[x_1, \dots, x_d] \otimes_{k[e_1, \dots, e_d]} k[[e_1, \dots, e_d]]$, and taking $S_d$-invariants --- again a kernel,
    preserved by the flat base change $k[e_1, \dots, e_d] \to k[[e_1, \dots, e_d]]$ --- yields $k[[e_1, \dots, e_d]]$.
    Under these identifications the maximal ideal of $\hat{A}$ is $(e_1, \dots, e_d)$, and $c$ is flat as the composition of
    the flat morphisms $\Spec(\hat{A}) \to \Spec(A) \to C^{(d)}$.

    It remains to check that $c^{-1}(C^{(d)} \setminus U^{(d)}) = V(e_d)$ as sets. We have
    $C^{(d)} \setminus U^{(d)} = \bigcup_{Q \in \mathrm{Supp}(\mdls)} Z_Q$ with $Z_Q = \{D : Q \in \mathrm{Supp}(D)\}$, as in
    the proof of \Cref{lemma:torsor_unramified_codim_one}. For $Q \neq P$ the divisor $Z_Q$ does not contain the point
    $Z_\ndls = dP$, so its ideal is not contained in $\mathfrak{m}_A$, and $c^{-1}(Z_Q) = \emptyset$. For $Q = P$, pull back
    along the finite surjective morphism $\Spec(\hat{B}) \to \Spec(\hat{A})$: the preimage of $Z_P$ in $C^d$ is
    $\bigcup_i p_i^{-1}(P)$, whose trace on $\Spec(\hat{B}) = \Spec(k[[x_1, \dots, x_d]])$ is
    $V(x_1) \cup \dots \cup V(x_d) = V(x_1 \cdots x_d) = V(e_d)$; by surjectivity, $c^{-1}(Z_P) = V(e_d)$ as claimed. Hence
    $c^{-1}(U^{(d)}) = V$, and likewise $\Spec(\hat{B}) \times_{C^d} U^d = W$.

    \textbf{(2)} Blowing up commutes with flat base change, and the center $Z_\ndls$ pulls back under $c$ to the closed
    point $V(e_1, \dots, e_d)$; hence
    $$\hat{Y} := \Bl_{(e_1, \dots, e_d)}\left(\Spec k[[e_1, \dots, e_d]]\right) \cong X_\ndls \times_{C^{(d)}} \Spec(\hat{A}),$$
    and the exceptional divisor $\hat{E} \subset \hat{Y}$ is the base change of $E_\ndls$ along
    $\Spec(\hat{A}/\hat{\mathfrak{m}}) = \Spec(\kappa(Z_\ndls))$, i.e.\ the projection $\mathrm{pr}: \hat{Y} \to X_\ndls$
    restricts to an isomorphism $\hat{E} \cong E_\ndls \cong \bP_k^{d-1}$.
    Let $\hat{\eta}$ be the generic point of $\hat{E}$. Then $\mathrm{pr}(\hat\eta) = \eta_\ndls$, and the induced local
    homomorphism of discrete valuation rings $\Oo_{X_\ndls, \eta_\ndls} \to \Oo_{\hat{Y}, \hat{\eta}}$ carries a uniformizer
    to a uniformizer (the ideal of $E_\ndls$ pulls back to the ideal of $\hat{E}$) and induces an isomorphism of residue
    fields (both are the function field of $\bP_k^{d-1}$); hence it induces an isomorphism on completions.
    The completion of $\Oo_{\hat{Y}, \hat{\eta}}$ is computed exactly as in
    \Cref{prop:blowup_local_ring_at_exceptional} (in the coordinates $e_1, \dots, e_d$, on the chart $u_d \neq 0$):
    we have $e_i = \frac{u_i}{u_d} e_d$, the local ring at
    $\hat\eta$ is $\left(\hat{A}[\frac{u_1}{u_d}, \dots, \frac{u_{d-1}}{u_d}]\right)_{(e_d)}$, its residue field is
    $\kappa = k(\frac{u_1}{u_d}, \dots, \frac{u_{d-1}}{u_d}) = k(\frac{e_1}{e_d}, \dots, \frac{e_{d-1}}{e_d})$, and its
    completion is $\kappa[[e_d]]$.
    Finally, $e_d$ is invertible in $\widehat{K}_{\eta_\ndls} = \kappa((e_d))$, so the composite
    $\Spec(\widehat{K}_{\eta_\ndls}) \to \Spec(\Oo_{\hat{Y}, \hat\eta}) \to \hat{Y} \to \Spec(\hat{A})$ factors through $V$,
    giving the asserted factorization of $\Spec(\widehat{K}_{\eta_\ndls}) \to U^{(d)}$ through $c_V$.

    \textbf{(3)} Let $h : \cP^{\times_k d} \to U^{(d)}$ denote the canonical morphism; it is finite, being the composition
    of the finite étale morphism $\cP^{\times_k d} \to U^d$ with the finite morphism $U^d \to U^{(d)}$. By
    \Cref{main:cor:sym-power-quotient},
    $\cP^{[d]} = \Xi \backslash \cP^{\times_k d} = \underline{\Spec}_{U^{(d)}}\left((h_* \Oo_{\cP^{\times_k d}})^{\Xi}\right)$.
    Pushforward along the affine morphism $h$ commutes with the flat base change $c_V : V \to U^{(d)}$, and the formation of
    $\Xi$-invariants is a kernel, hence also commutes with flat base change; therefore
    $$ \cP^{[d]} \times_{U^{(d)}} V \cong \Xi \backslash \left(\cP^{\times_k d} \times_{U^{(d)}} V\right).$$
    By step (1), $U^d \times_{U^{(d)}} V \cong W$, so
    $\cP^{\times_k d} \times_{U^{(d)}} V \cong \cP^{\times_k d} \times_{U^d} W$.
    For each $i$, the composition $W \to \Spec(k[[x_1, \dots, x_d]]) \to \Spec(k[[x]]) \to C$ --- the middle map being the
    completion of the projection $p_i : C^d \to C$, i.e.\ $x \mapsto x_i$ --- factors through $\Spec(k((x))) \to U$, since
    $x_i$ is invertible on $W$; this composite is precisely $q_i$, and it commutes with $p_i : U^d \to U$. Since
    $\cP^{\times_k d} = p_1^{-1}\cP \times_{U^d} \cdots \times_{U^d} p_d^{-1}\cP$, base changing to $W$ gives an isomorphism
    $$ \cP^{\times_k d} \times_{U^d} W \cong q_1^{-1}\cP_x \times_W \cdots \times_W q_d^{-1}\cP_x, $$
    which is equivariant for the actions of $G^d$ (in particular of $\KG$) and of $S_d$ --- the latter acting on $W$ by
    permuting the $x_i$, compatibly with $q_{\sigma(i)} = q_i \circ \sigma$. Passing to $\Xi$-quotients yields the displayed
    isomorphism.

    For the final statement: the schemes $V$ and $W$, the maps $q_i$ and $\hat{r}$, the $\Xi$-action, and the point
    $\Spec(\kappa((e_d))) \to V$ of (2) are all constructed from $d$ alone; by (3), the torsor
    $\cP^{[d]}|_{\Spec(\widehat{K}_{\eta_\ndls})}$ is the pullback along $\Spec(\kappa((e_d))) \to V$ of a torsor constructed
    from $\cP_x$ and these data. Since the ramification at $\eta_\ndls$ is by definition read off from
    $\cP^{[d]}|_{\Spec(\widehat{K}_{\eta_\ndls})}$ together with the discrete valuation ring $\kappa[[e_d]]$
    (\Cref{definition:tamely_ramified_in_codim_1}, \Cref{definition:local_ramification_boundness}), it depends only on
    $(\cP_x, d)$.
\end{proof}

\begin{cor}[Reduction to $\bP^1$]\label{cor:reduction_to_P1}
    Let $C, C'$ be smooth projective geometrically connected curves over the algebraically closed field $k$, let
    $\cP \to U = C \setminus \mdls$ and $\cP' \to U' = C' \setminus \mdls'$ be $G$-torsors, let
    $P \in \mathrm{Supp}(\mdls)(k)$ and $P' \in \mathrm{Supp}(\mdls')(k)$, and let $x, x'$ be uniformizers at $P, P'$
    respectively. Assume that $\cP_x \cong \cP'_{x'}$ as $G$-torsors, compatibly with the isomorphism
    $k((x)) \cong k((x'))$, $x \mapsto x'$.
    Then for every $d \geq 1$ and every $r$, the ramification of $\cP^{[d]}$ at $\eta_{dP}$ is bounded by $r$ if and only if
    the ramification of $\cP'^{[d]}$ at $\eta_{dP'}$ is; in particular, the one is unramified (resp.\ tamely ramified) at
    $\eta_{dP}$ if and only if the other is at $\eta_{dP'}$.
\end{cor}
\begin{proof}
    Immediate from \Cref{lemma:formal_local_invariance}: both $\cP^{[d]}|_{\Spec(\widehat{K}_{\eta_{dP}})}$ and
    $\cP'^{[d]}|_{\Spec(\widehat{K}_{\eta_{dP'}})}$ are identified, compatibly with the discrete valuation ring
    $\kappa[[e_d]]$, with the pullback along $\Spec(\kappa((e_d))) \to V$ of
    $\Xi \backslash (q_1^{-1}\cP_x \times_W \cdots \times_W q_d^{-1}\cP_x)$.
\end{proof}

\subsection*{The Invariant-Theory Scaffold}

All three computations of \Cref{section:symmetric_power_computations} identify
the symmetric-power cover at $\eta_\ndls$ by the same invariant-theoretic
mechanism, which we set up once. Let $G$ be a finite abelian group and let
$\cP = \Spec S \to \bG_m = \Spec R$, $R = k[x, x^{-1}]$, be an affine
$G$-torsor. Write
\[
    R^{\otimes_k d} = k[x_1^{\pm 1}, \dots, x_d^{\pm 1}],
\]
the coordinate ring of $U^d = \bG_m^d$, and let
$S^{\otimes_k d}$ be the coordinate ring of the $G^d$-torsor
$\cP^{\times_k d} = p_1^{-1}\cP \times_{U^d} \cdots \times_{U^d} p_d^{-1}\cP$
over $U^d$.

The $d$-fold product torsor
$p_1^{-1}\cP \otimes^G \cdots \otimes^G p_d^{-1}\cP$ on $U^d$ --- the
contracted product along the sum map $G^d \to G$ --- is the quotient of
$\cP^{\times_k d}$ by the kernel
$\KG = \{(g_1, \dots, g_d) \in G^d \mid g_1 + \dots + g_d = 0\}$ of the sum
map. We identify $\KG \cong G^{d-1}$ via
$(g_1, \dots, g_{d-1}) \mapsto (g_1,\, g_2 - g_1,\, \dots,\, -g_{d-1})$; its
coordinate ring is the invariant ring $\left(S^{\otimes_k d}\right)^{G^{d-1}}$,
and the residual action of $G = G^d / \KG$ is induced by
$(g, 0, \dots, 0) \in G^d$.
The symmetric group $S_d$ acts on everything by permuting the $d$ factors,
and by \Cref{main:cor:sym-power-quotient} the symmetric-power torsor is the
$\Xi = \KG \rtimes S_d$ quotient:
\[
    \cP^{[d]}\big|_{U^{(d)}}
    = \Spec \left( \left( S^{\otimes_k d} \right)^{\Xi} \right)
    = \Spec \left( \left( \left( S^{\otimes_k d} \right)^{G^{d-1}} \right)^{S_d} \right).
\]
On the base, the $S_d$-invariants form the ring of symmetric Laurent
polynomials
\[
    \left(R^{\otimes_k d}\right)^{S_d} = k[e_1, e_2, \dots, e_d, e_d^{-1}],
\]
where $e_1, \dots, e_d$ are the elementary symmetric polynomials in
$x_1, \dots, x_d$ (a Laurent polynomial is symmetric if and only if, after
clearing the denominator $e_d^N = (x_1 \cdots x_d)^N$, its numerator is a
symmetric polynomial). Finally, by
\Cref{lemma:formal_local_invariance}(2)--(3) applied to $C = \bP^1$,
$\mdls = d \cdot 0 + \infty$, the restriction of $\cP^{[d]}$ to the complete
local field $\widehat{K}_\eta = \kappa(\eta)((e_d))$ of
\Cref{eq:complete_field_at_generic_point} is the base change of
$\Spec\left( \left( S^{\otimes_k d} \right)^{\Xi} \right)$ along
$k[e_1, \dots, e_d, e_d^{-1}] \to \widehat{K}_\eta$. In each of the three
computations, the task is therefore to identify the invariant ring
$\left( S^{\otimes_k d} \right)^{\Xi}$ explicitly as the algebra of an
equation of the same shape as that of $\cP$, with datum a \emph{symmetric}
function of the $x_i$, and then to bound the valuation of that datum in
$\widehat{K}_\eta$.

\subsection*{Compatibilities of the Symmetric Power}

% --- NEW (not in the current chapter): the symmetric power commutes with
%     fiber products and with induced torsors. These are the two formal
%     compatibilities that the assembly of the general abelian case (2.5.4)
%     and the tame case (2.5.1) consume. Both proofs follow the same pattern
%     as prop:witt_cover_identification: produce an equivariant morphism of
%     torsors and conclude by prop:torsor_morphism_iso. ---

Recall from \Cref{main:cor:sym-power-quotient} that
$\cP^{[d]} = \Xi_G \backslash \cP^{\times_k d}$, where
$\Xi_G = \KG \rtimes S_d$. The symmetric power is compatible with the two
standard operations on the structure group: fiber products and induction
along injections.

\begin{lemma}[symmetric powers and fiber products]\label{lemma:sym_power_product}
    Let $G_1, G_2$ be finite abelian groups, let $\cP_1, \cP_2$ be a $G_1$-,
    resp.\ $G_2$-torsor on $U_{\text{ét}}$, and let
    $\cP = \cP_1 \times_U \cP_2$ be the associated $(G_1 \times G_2)$-torsor.
    Then there is a canonical isomorphism of $(G_1 \times G_2)$-torsors over
    $U^{(d)}$
    \[
        \cP^{[d]} \;\cong\; \cP_1^{[d]} \times_{U^{(d)}} \cP_2^{[d]}.
    \]
\end{lemma}
\begin{proof}
    Write $G = G_1 \times G_2$. We have
    $\cP^{\times_k d} = \cP_1^{\times_k d} \times_{U^d} \cP_2^{\times_k d}$,
    and under this identification
    $K_G = K_{G_1} \times K_{G_2}$ inside $G^d = G_1^d \times G_2^d$, with
    $S_d$ acting diagonally. Consider the composition
    \[
        \cP^{\times_k d}
        \longrightarrow \cP_1^{\times_k d} \times_{U^d} \cP_2^{\times_k d}
        \longrightarrow
        \bigl(\Xi_{G_1} \backslash \cP_1^{\times_k d}\bigr)
        \times_{U^{(d)}}
        \bigl(\Xi_{G_2} \backslash \cP_2^{\times_k d}\bigr)
        = \cP_1^{[d]} \times_{U^{(d)}} \cP_2^{[d]}.
    \]
    The first projection $\cP^{\times_k d} \to \cP_1^{[d]}$ is invariant under
    $K_{G_1} \rtimes S_d$ by construction and under $K_{G_2}$ because
    $K_{G_2}$ acts only on the second factor; symmetrically for the second
    projection. Hence the composition is $\Xi_G$-invariant and descends to a
    morphism
    \[
        \phi \colon \cP^{[d]} = \Xi_G \backslash \cP^{\times_k d}
        \longrightarrow \cP_1^{[d]} \times_{U^{(d)}} \cP_2^{[d]}
    \]
    over $U^{(d)}$. The residual $G$-action on the source is induced by
    $(g_1, g_2, 0, \dots, 0) \in G^d$, which maps under the projections to the
    residual $G_1$- and $G_2$-actions on the factors of the target; so $\phi$
    is $G$-equivariant. The source is a $G$-torsor
    (\Cref{main:cor:sym-power-quotient}), and the target is a fiber product of a
    $G_1$-torsor and a $G_2$-torsor, hence a $G$-torsor. A $G$-equivariant
    morphism of $G$-torsors is an isomorphism
    (\Cref{main:prop:torsor_morphism_iso}).
\end{proof}

\begin{lemma}[symmetric powers and induced torsors]\label{lemma:sym_power_induced}
    Let $\iota : H \hookrightarrow G$ be an injective homomorphism of finite
    abelian groups and let $\cQ$ be an $H$-torsor on $U_{\text{ét}}$. Then
    there is a canonical isomorphism of $G$-torsors over $U^{(d)}$
    \[
        (\iota_* \cQ)^{[d]} \;\cong\; \iota_*\bigl(\cQ^{[d]}\bigr).
    \]
\end{lemma}
\begin{proof}
    The canonical map $j \colon \cQ \to \iota_*\cQ = (\cQ \times G)/H$,
    $q \mapsto [(q, 1)]$, satisfies $j(qh) = j(q)\cdot\iota(h)$. Its $d$-fold
    product $\cQ^{\times_k d} \to (\iota_*\cQ)^{\times_k d}$ is equivariant
    for $H^d \to G^d$ (which carries $K_H$ into $K_G$) and commutes with the
    $S_d$-actions, hence descends to an $H$-equivariant morphism of the
    quotients
    \[
        \psi \colon \cQ^{[d]} = \Xi_H \backslash \cQ^{\times_k d}
        \longrightarrow \Xi_G \backslash (\iota_*\cQ)^{\times_k d}
        = (\iota_*\cQ)^{[d]},
    \]
    where the residual $H = H^d/K_H$ acts on the target through
    $\iota : H \to G = G^d/K_G$. By the universal property of induction, an
    $H$-equivariant morphism from the $H$-torsor $\cQ^{[d]}$ to (the
    restriction of) the $G$-torsor $(\iota_*\cQ)^{[d]}$ extends uniquely to a
    $G$-equivariant morphism
    \[
        \iota_*\bigl(\cQ^{[d]}\bigr) \longrightarrow (\iota_*\cQ)^{[d]},
        \qquad [(z, g)] \longmapsto \psi(z) \cdot g,
    \]
    which is well defined precisely because $\psi$ is $H$-equivariant. A
    $G$-equivariant morphism of $G$-torsors over $U^{(d)}$ is an isomorphism
    (\Cref{main:prop:torsor_morphism_iso}).
\end{proof}

\subsection*{A Symmetric Regularity Lemma}

The valuation bound is in all cases an instance of the following lemma, applied
with $y_i := 1/x_i$. Note that the elementary symmetric polynomials in the
$y_i$ are related to those in the $x_i$ by
\begin{equation}\label{eq:elementary_symmetric_inverse_variables}
    e_k(y_1, \dots, y_d)
    = \sum_{1 \le i_1 < \dots < i_k \le d} \frac{1}{x_{i_1} \cdots x_{i_k}}
    = \frac{e_{d-k}(x_1, \dots, x_d)}{e_d(x_1, \dots, x_d)},
\end{equation}
which for $1 \leq k \leq d-1$ lies in the residue field
$\kappa(\eta) = k(e_1/e_d, \dots, e_{d-1}/e_d)$ of
\Cref{eq:complete_field_at_generic_point}.

\begin{lemma}[symmetric regularity]\label{lemma:symmetric_regularity}
    Let $\Phi \in k[y_1, \dots, y_d]^{S_d}$ be symmetric with every monomial of total
    degree $< d$. Then $\Phi \in \kappa(\eta) = k(e_1/e_d, \dots, e_{d-1}/e_d)$; in
    particular $\Phi$ is regular at $\eta_\ndls$.
\end{lemma}

\begin{proof}
    $\Phi$ is a $k$-combination of monomial symmetric functions $m_\lambda(y)$ with
    $|\lambda| < d$. Each $m_\lambda(y)$ is an isobaric polynomial in the elementary
    symmetric functions $e_\mu(y) = \prod_i e_{\mu_i}(y)$ with $\mu \vdash |\lambda|$;
    every part $\mu_i \leq |\lambda| < d$, so only $e_1(y), \dots, e_{d-1}(y)$ occur
    --- never $e_d(y)$. By \Cref{eq:elementary_symmetric_inverse_variables}, each
    $e_k(y)$ with $1 \leq k \leq d-1$ lies in $\kappa(\eta)$. Hence
    $\Phi \in \kappa(\eta) \subset \widehat\Oo_\eta = \kappa(\eta)[[e_d]]$.
\end{proof}

% ---------------------------------------------------------------------------
%  2.5  The Symmetric-Power Computations
%  [Sources: theorem:ramification_after_blowup_P1 and lemma:realization_over_Gm
%   of RamificationAfterBlowupIsTame.tex, split into per-case subsections with
%   a common skeleton; RamificationAfterBlowupPPower.tex for the p^r case;
%   subsection 2.5.4 (general finite abelian G) is NEW --- it replaces the
%   one-line stub "Idea: decompose G by the structure theorem".]
% ---------------------------------------------------------------------------

\section{The Symmetric-Power Computations}\label{section:symmetric_power_computations}

We now prove \Cref{theorem:torsors_after_blowup_is_tame}. The proof proceeds
through the three cyclic cases --- $G = \Z/e\Z$ with $\gcd(e,p) = 1$,
$G = \Z/p\Z$, and $G = \Z/p^r\Z$ --- each treated by the \emph{same}
three-step computation, one column of the table of
\Cref{section:cyclic_extension_theories} at a time:
\begin{enumerate}
    \item \textbf{Realize.} The restriction $\cP_x$ of $\cP$ to the punctured
    formal disc at $P$ is realized, using the reduced representatives of
    \Cref{section:cyclic_extension_theories}, by an explicit equation over
    $\bG_m \subset \bP^1$; by \Cref{cor:reduction_to_P1} we may then assume
    $C = \bP^1_k$, $\ndls = d \cdot 0$, $U = \bG_m$, and $\cP$ is this
    explicit torsor.
    \item \textbf{Identify.} Using the invariant-theory scaffold of
    \Cref{section:blowup_prelude}, the restriction of $\cP^{[d]}$ to the
    complete local field $\widehat K_\eta = \kappa(\eta)((e_d))$ of
    \Cref{eq:complete_field_at_generic_point} is identified as the cover
    defined by an equation \emph{of the same shape}, whose datum is the
    symmetrization of the original datum: the product $\prod_i f(x_i)$ in the
    Kummer case, the sum $\sum_i f(x_i)$ in the Artin--Schreier case, and the
    Witt sum $\bigoplus_i \vec f(x_i)$ in the Artin--Schreier--Witt case.
    \item \textbf{Count.} The valuation of the symmetrized datum in
    $\widehat K_\eta$ is bounded --- using
    \Cref{lemma:symmetric_regularity} and the pole bounds built into the
    reduced representative --- and the ramification criterion of the relevant
    theory yields the conclusion: \emph{tame} in the Kummer case,
    \emph{unramified} in the two wild cases.
\end{enumerate}
Finally, in \Cref{subsection:general_abelian_case}, the general finite
abelian case is assembled from the three cyclic cases by decomposing $G$ into
primary cyclic factors, using the compatibility lemmas of
\Cref{section:blowup_prelude}.

Throughout this section $k$ is algebraically closed
(\Cref{rem:algebraically_closed_harmless}), $P \in C(k)$, $\ndls = nP$, and
$d = \deg \ndls = n$.

% ---------------------------------------------------------------------
\subsection{The Tame Case}\label{subsection:tame_case}
% ---------------------------------------------------------------------

\subsubsection*{Step 1: realization over $\bG_m$}

\begin{lemma}[Realization over $\bG_m$, Kummer case]\label{lemma:realization_over_Gm_tame}
    Let $k$ be algebraically closed of characteristic $p > 0$, and let $\cQ$ be a
    tamely ramified $G$-torsor over $\Spec(k((x)))$. Then there exist an integer
    $e$ prime to $p$, an injective homomorphism $\iota : \Z/e\Z \hookrightarrow G$,
    and an integer $a$ prime to $e$, such that
    $$\cP'_0 := \Spec\left(k[x, x^{-1}][T]/(T^e - x^a)\right)$$
    is a $\Z/e\Z$-torsor over $\bG_m$, tamely ramified at $0$ and $\infty$, and the
    induced $G$-torsor $\cP' := \iota_*\cP'_0$ satisfies $\cP'_x \cong \cQ$.
    Here $T^e = x^a$ is regarded as a $\Z/e\Z$-torsor via a fixed primitive $e$-th
    root of unity $\zeta_e \in k$.
\end{lemma}
\begin{proof}
    By \Cref{lemma:kummer_reduced_representatives}(1), the tame quotient $G^t$ of
    $\Gal(k((x))^{sep}/k((x)))$ is procyclic; fix a topological generator $\gamma$.
    Let $\rho: \Gal(k((x))^{sep}/k((x))) \to G$ be the continuous homomorphism
    classifying $\cQ$. Tameness means that $\rho$ kills the wild inertia subgroup,
    so $\rho$ factors through $G^t$ and is determined by $g := \rho(\gamma)$, whose
    order $e$ is prime to $p$. Let $\iota: \Z/e\Z \xrightarrow{\sim} \langle g \rangle \subset G$
    be given by $1 \mapsto g$, so that $\rho = \iota \circ \chi$ for a surjective
    character $\chi : G^t \to \Z/e\Z$.
    By \Cref{lemma:kummer_reduced_representatives}(2), the classes of $x^a$,
    $a \in \Z/e\Z$, exhaust $\mathrm{Hom}_{\mathrm{cont.}}(G^t, \Z/e\Z) \cong \Z/e\Z$, the class
    of $x^a$ being surjective precisely when $\gcd(a, e) = 1$. Choose $a$ with the
    class of $x^a$ equal to $\chi$; then $\cQ \cong \iota_*(T^e = x^a)$ over
    $k((x))$. The same equation over $k[x, x^{-1}]$ defines a $\Z/e\Z$-torsor
    $\cP'_0$ over $\bG_m$ (étale since $e$ is invertible in $k$), with
    $(\iota_*\cP'_0)_x \cong \cQ$; it is tamely ramified at $0$ and at $\infty$
    since its local characters there factor through the tame quotients.
\end{proof}

\subsubsection*{Step 2: identification of the symmetric-power cover}

\begin{prop}[identification of the product torsor, Kummer case]\label{prop:kummer_cover_identification}
    Let $G = \Z/e\Z$ with $e$ prime to $p$, let $f \in R^\times$ be a unit of
    $R = k[x, x^{-1}]$, and let $\cP' = \Spec S \to \bG_m$, $S = R[X]/(X^e - f)$, be
    the associated $G$-torsor. In the notation of the scaffold of
    \Cref{section:blowup_prelude}, set
    \[
        Y := X_1 X_2 \cdots X_d \in S^{\otimes_k d},
        \qquad
        F := \prod_{i=1}^{d} f(x_i) \in R^{\otimes_k d}.
    \]
    Then:
    \begin{enumerate}
        \item \label{item:kummer_cover_invariance} $Y$ is $\KG$-invariant, and
        \[
            \bigl(S^{\otimes_k d}\bigr)^{G^{d-1}}
            \;\cong\;
            R^{\otimes_k d}[Y]\big/\bigl(Y^e - F\bigr)
        \]
        as $G$-torsors over $R^{\otimes_k d}$ (the residual $G$ acting by
        $Y \mapsto \zeta_e^{\,g}\, Y$).
        \item \label{item:kummer_cover_symmetry} $F$ is a symmetric Laurent
        polynomial, i.e.\ lies in
        $\bigl(R^{\otimes_k d}\bigr)^{S_d} = k[e_1, \dots, e_d, e_d^{-1}]$, and
        \[
            \Bigl(\bigl(S^{\otimes_k d}\bigr)^{G^{d-1}}\Bigr)^{S_d}
            \;\cong\;
            k[e_1, \dots, e_d, e_d^{-1}][Y]\big/\bigl(Y^e - F\bigr).
        \]
    \end{enumerate}
    Consequently, restricting to the complete local field
    $\widehat K_\eta = \kappa(\eta)((e_d))$ of
    \Cref{eq:complete_field_at_generic_point},
    \[
        \cP'^{[d]}\big|_{\Spec \widehat K_\eta}
        \;\cong\;
        \Spec\; \widehat K_\eta[Z]\big/\bigl(Z^e - F\bigr).
    \]
\end{prop}
\begin{proof}
    \ref{item:kummer_cover_invariance}
    Recall that $g \in G = \Z/e\Z$ acts on $X$ by $g(X) = \zeta_e^{\,g} X$. Hence
    $(g_1, \dots, g_{d-1}) \in \KG \cong G^{d-1}$ acts on the generators by
    \[
        X_1 \mapsto \zeta_e^{\,g_1} X_1, \qquad
        X_i \mapsto \zeta_e^{\,g_i - g_{i-1}} X_i \ (1 < i < d), \qquad
        X_d \mapsto \zeta_e^{\,-g_{d-1}} X_d,
    \]
    so $Y = X_1 \cdots X_d$ transforms by the telescoping product
    $\zeta_e^{\,g_1 + (g_2 - g_1) + \dots - g_{d-1}} = 1$: it is invariant.
    Moreover $Y^e = \prod_i X_i^e = \prod_i f(x_i) = F$.
    Since $F$ is a unit of $R^{\otimes_k d}$ and $e$ is invertible in $k$,
    $T := R^{\otimes_k d}[Y]/(Y^e - F)$ is a $\Z/e\Z$-torsor over $U^d$ (finite
    étale, with $G$ acting by $Y \mapsto \zeta_e^{\,g} Y$), and
    $Y \mapsto X_1 \cdots X_d$ defines an $R^{\otimes_k d}$-algebra homomorphism
    $\phi : T \to \bigl(S^{\otimes_k d}\bigr)^{G^{d-1}}$. The target is the
    $d$-fold product torsor (the contracted product of the $p_i^{-1}\cP'$ along
    the sum map $G^d \to G$), a $G$-torsor over $U^d$ by
    \Cref{main:prop:symmetric_power_torsor} and the surrounding discussion; its
    residual $G$-action is induced by $(g, 0, \dots, 0) \in G^d$, under which
    $Y \mapsto \zeta_e^{\,g} Y$. So $\Spec \phi$ is a $G$-equivariant morphism of
    $G$-torsors over $U^d$, hence an isomorphism by
    \Cref{main:prop:torsor_morphism_iso}.

    \ref{item:kummer_cover_symmetry} $S_d$ permutes the $x_i$ and the $X_i$, so it
    fixes $Y$ and $F$; in particular
    $F \in \bigl(R^{\otimes_k d}\bigr)^{S_d} = k[e_1, \dots, e_d, e_d^{-1}]$ (as
    recorded in the scaffold). $T$ is free over $R^{\otimes_k d}$ with basis the
    $S_d$-invariant monomials $Y^j$, $0 \leq j < e$; writing an element in this
    basis, it is $S_d$-invariant if and only if all its coefficients are, whence
    \[
        T^{S_d}
        = \bigl(R^{\otimes_k d}\bigr)^{S_d}[Y]\big/\bigl(Y^e - F\bigr)
        = k[e_1, \dots, e_d, e_d^{-1}][Y]\big/\bigl(Y^e - F\bigr).
    \]
    The final statement follows from the scaffold: by
    \Cref{main:cor:sym-power-quotient}, $\cP'^{[d]}$ is the $\Xi$-quotient computed
    above, and its restriction to $\Spec \widehat K_\eta$ along
    $k[e_1, \dots, e_d, e_d^{-1}] \to \widehat K_\eta$
    (\Cref{lemma:formal_local_invariance}) is the Kummer cover of the image of
    $F$.
\end{proof}

\subsubsection*{Step 3: the valuation count and the conclusion}

\begin{theorem}[the tame case]\label{theorem:tame_after_blowup}
    Let $G$ be a finite abelian group and let $\cP \to U = C \setminus \mdls$ be a
    $G$-torsor which is tamely ramified at $P$. Then $\cP^{[d]}$ is tamely
    ramified at $\eta_\ndls$.
\end{theorem}
\begin{proof}
    \emph{Realize.} The restriction $\cQ = \cP_x$ is tamely ramified, so
    \Cref{lemma:realization_over_Gm_tame} produces $e$ prime to $p$,
    $\iota: \Z/e\Z \hookrightarrow G$, $a$ prime to $e$, and the torsor
    $\cP'_0 : T^e = x^a$ over $\bG_m$ with $(\iota_*\cP'_0)_x \cong \cP_x$. By
    \Cref{cor:reduction_to_P1}, the ramification of $\cP^{[d]}$ at $\eta_\ndls$
    agrees with that of $(\iota_*\cP'_0)^{[d]}$ at $\eta_{d \cdot 0}$. By
    \Cref{lemma:sym_power_induced},
    $(\iota_*\cP'_0)^{[d]} \cong \iota_*\bigl(\cP'^{[d]}_0\bigr)$, and by
    \Cref{lemma:ramification_of_induced_torsor} it is tamely ramified at
    $\eta_{d \cdot 0}$ if and only if $\cP'^{[d]}_0$ is. So we may assume
    $G = \Z/e\Z$ and $\cP = \cP'_0$.

    \emph{Identify.} \Cref{prop:kummer_cover_identification}, with $f = x^a$,
    identifies
    \[
        \cP'^{[d]}_0\big|_{\Spec \widehat K_\eta}
        \;\cong\;
        \Spec\; \widehat K_\eta[Z]\big/\bigl(Z^e - F\bigr),
        \qquad
        F = \prod_{i=1}^{d} x_i^a = (x_1 \cdots x_d)^a = e_d^{\,a}.
    \]

    \emph{Count.} In $\widehat K_\eta = \kappa(\eta)((e_d))$ we have
    $v_{\widehat K_\eta}(F) = v_{\widehat K_\eta}(e_d^{\,a}) = a$. By
    \Cref{theorem:kummer_ramification} (all of whose cases are tamely ramified,
    since $\gcd(e, p) = 1$), the extension defined by $Z^e = e_d^{\,a}$ is tamely
    ramified --- explicitly, its ramification index is $e/\gcd(a, e)$, prime to
    $p$. Hence $\cP'^{[d]}_0$, and therefore $\cP^{[d]}$, is tamely ramified at
    $\eta_\ndls$.
\end{proof}

% ---------------------------------------------------------------------
\subsection{\texorpdfstring{The Wild Case $G = \Z/p\Z$}{The Wild Case G = Z/pZ}}\label{subsection:wild_case_Zp}
% ---------------------------------------------------------------------

\subsubsection*{Step 1: realization over $\bG_m$}

\begin{lemma}[Realization over $\bG_m$, Artin--Schreier case]\label{lemma:realization_over_Gm_AS}
    Let $k$ be algebraically closed of characteristic $p > 0$, let $G = \Z/p\Z$,
    and let $\cQ$ be a $G$-torsor over $\Spec(k((x)))$ with ramification bounded
    by $d$ (\Cref{definition:local_ramification_boundness}). Then there exists
    $f \in x^{-1}k[x^{-1}]$, either $f = 0$ or
    $f = cx^{-m} + a_{-m+1}x^{-m+1} + \dots + a_{-1}x^{-1}$ with $c \neq 0$,
    $p \nmid m$ and $m < d$, such that
    $$\cP' := \Spec\left(k[x, x^{-1}][X]/(X^p - X - f)\right)$$
    is a $\Z/p\Z$-torsor over $\bG_m$ with $\cP'_x \cong \cQ$. Moreover $\cP'$
    extends to a torsor over $\bP^1_k \setminus \{0\}$ --- in particular it is
    unramified at $\infty$ --- and it has ramification bounded by $d$ at $0$.
\end{lemma}
\begin{proof}
    If $\cQ$ is trivial, take $f = 0$. Otherwise, since $\Z/p\Z$ is simple and
    $k((x))$ admits no non-trivial unramified extensions, $\cQ = \Spec(F)$ for a
    cyclic extension $F/k((x))$ of degree $p$ together with an isomorphism
    $\Gal(F/k((x))) \cong \Z/p\Z$. By Artin--Schreier theory
    (\Cref{theorem:artin_schreier_ramification}) such data correspond to non-zero
    classes in $k((x))/\wp(k((x)))$, and by
    \Cref{lemma:as_reduced_representatives} the class of $\cQ$ admits a reduced
    representative $f \in x^{-1}k[x^{-1}]$, non-zero with pole order $m$ prime to
    $p$. The unique ramification break of $F = k((x))[X]/(X^p - X - f)$ is at
    $m$ (\Cref{theorem:artin_schreier_ramification}(3)), so the ramification of
    $\cQ$ is bounded by $d$ if and only if $m < d$.
    Finally, $f \in k[x, x^{-1}]$, so the same equation defines a finite étale
    $\Z/p\Z$-cover $\cP' \to \bG_m$ (a torsor via $X \mapsto X + 1$) with
    $\cP'_x \cong \cQ$; and writing $y = x^{-1}$, we have $f \in y\,k[y]$, so
    $X^p - X - f$ defines a finite étale cover of
    $\mathbb{A}^1_k = \Spec(k[y]) = \bP^1_k \setminus \{0\}$ extending $\cP'$,
    whence $\cP'$ is unramified at $\infty$.
\end{proof}

\subsubsection*{Step 2: identification of the symmetric-power cover}

\begin{prop}[identification of the product torsor, Artin--Schreier case]\label{prop:as_cover_identification}
    Let $G = \Z/p\Z$, let $f \in R = k[x, x^{-1}]$, and let
    $\cP' = \Spec S \to \bG_m$, $S = R[X]/(X^p - X - f)$, be the associated
    $G$-torsor. In the notation of the scaffold, set
    \[
        Y := X_1 + X_2 + \dots + X_d \in S^{\otimes_k d},
        \qquad
        \alpha := \sum_{i=1}^{d} f(x_i) \in R^{\otimes_k d}.
    \]
    Then:
    \begin{enumerate}
        \item \label{item:as_cover_invariance} $Y$ is $\KG$-invariant, and
        \[
            \bigl(S^{\otimes_k d}\bigr)^{G^{d-1}}
            \;\cong\;
            R^{\otimes_k d}[Y]\big/\bigl(Y^p - Y - \alpha\bigr)
        \]
        as $G$-torsors over $R^{\otimes_k d}$ (the residual $G$ acting by
        $Y \mapsto Y + g$).
        \item \label{item:as_cover_symmetry} $\alpha$ is a symmetric Laurent
        polynomial, i.e.\ lies in
        $\bigl(R^{\otimes_k d}\bigr)^{S_d} = k[e_1, \dots, e_d, e_d^{-1}]$, and
        \[
            \Bigl(\bigl(S^{\otimes_k d}\bigr)^{G^{d-1}}\Bigr)^{S_d}
            \;\cong\;
            k[e_1, \dots, e_d, e_d^{-1}][Y]\big/\bigl(Y^p - Y - \alpha\bigr).
        \]
    \end{enumerate}
    Consequently, restricting to the complete local field
    $\widehat K_\eta = \kappa(\eta)((e_d))$ of
    \Cref{eq:complete_field_at_generic_point},
    \[
        \cP'^{[d]}\big|_{\Spec \widehat K_\eta}
        \;\cong\;
        \Spec\; \widehat K_\eta[Z]\big/\bigl(Z^p - Z - \alpha\bigr).
    \]
\end{prop}
\begin{proof}
    \ref{item:as_cover_invariance}
    Recall that $g \in G = \Z/p\Z$ acts on $X$ by $g(X) = X + g$. Hence
    $(g_1, \dots, g_{d-1}) \in \KG \cong G^{d-1}$ acts on the generators by
    \[
        X_1 \mapsto X_1 + g_1, \qquad
        X_i \mapsto X_i - g_{i-1} + g_i \ (1 < i < d), \qquad
        X_d \mapsto X_d - g_{d-1},
    \]
    so $Y = X_1 + \dots + X_d$ changes by the telescoping sum
    $g_1 + (g_2 - g_1) + \dots + (-g_{d-1}) = 0$: it is invariant. Moreover
    $\wp(Y) = \sum_i \wp(X_i) = \sum_i f(x_i) = \alpha$ ($\wp$ is additive).
    By \Cref{lemma:asw_etale_torsor} (with $r = 1$ and $A = R^{\otimes_k d}$),
    $T := R^{\otimes_k d}[Y]/(Y^p - Y - \alpha)$ is a $\Z/p\Z$-torsor over $U^d$,
    and $Y \mapsto X_1 + \dots + X_d$ defines an $R^{\otimes_k d}$-algebra
    homomorphism $\phi : T \to \bigl(S^{\otimes_k d}\bigr)^{G^{d-1}}$. The target
    is the $d$-fold product torsor, a $G$-torsor over $U^d$ by
    \Cref{main:prop:symmetric_power_torsor} and the surrounding discussion; its
    residual $G$-action is induced by $(g, 0, \dots, 0) \in G^d$, under which
    $Y \mapsto Y + g$. So $\Spec \phi$ is a $G$-equivariant morphism of
    $G$-torsors over $U^d$, hence an isomorphism by
    \Cref{main:prop:torsor_morphism_iso}.

    \ref{item:as_cover_symmetry} $S_d$ permutes the $x_i$ and the $X_i$, so it
    fixes $Y$ and $\alpha$; in particular
    $\alpha \in \bigl(R^{\otimes_k d}\bigr)^{S_d} = k[e_1, \dots, e_d, e_d^{-1}]$.
    $T$ is free over $R^{\otimes_k d}$ with basis the $S_d$-invariant monomials
    $Y^j$, $0 \leq j < p$ (\Cref{lemma:asw_etale_torsor}); writing an element in
    this basis, it is $S_d$-invariant if and only if all its coefficients are,
    whence
    \[
        T^{S_d}
        = k[e_1, \dots, e_d, e_d^{-1}][Y]\big/\bigl(Y^p - Y - \alpha\bigr).
    \]
    The final statement follows from the scaffold exactly as in
    \Cref{prop:kummer_cover_identification}.
\end{proof}

\subsubsection*{Step 3: the degree count and the conclusion}

\begin{theorem}[the wild case, $G = \Z/p\Z$]\label{theorem:wild_Zp_after_blowup}
    Let $G = \Z/p\Z$ and let $\cP \to U = C \setminus \mdls$ be a $G$-torsor with
    ramification bounded by $\ndls = nP$ at $P$. Then $\cP^{[d]}$ is unramified
    at $\eta_\ndls$.
\end{theorem}
\begin{proof}
    \emph{Realize.} The restriction $\cQ = \cP_x$ is a $\Z/p\Z$-torsor with
    ramification bounded by $d$, so \Cref{lemma:realization_over_Gm_AS} produces
    $\cP' : X^p - X = f$ over $\bG_m$ with $\cP'_x \cong \cP_x$, where $f$ is
    reduced with pole order $m < d$ (or $f = 0$). By \Cref{cor:reduction_to_P1},
    we may assume $\cP = \cP'$.

    \emph{Identify.} \Cref{prop:as_cover_identification} identifies
    \[
        \cP'^{[d]}\big|_{\Spec \widehat K_\eta}
        \;\cong\;
        \Spec\; \widehat K_\eta[Z]\big/\bigl(Z^p - Z - \alpha\bigr),
        \qquad
        \alpha = \sum_{i=1}^{d} f(x_i).
    \]

    \emph{Count.} In the variables $y_i = x_i^{-1}$, each $f(x_i)$ is a
    polynomial in $y_i$ of degree $m$ with zero constant term, so $\alpha$ is a
    symmetric polynomial in $y_1, \dots, y_d$ all of whose monomials have total
    degree $\leq m < d$. By \Cref{lemma:symmetric_regularity},
    $\alpha \in \kappa(\eta) \subset \widehat\Oo_\eta = \kappa(\eta)[[e_d]]$; in
    particular $v_{\widehat K_\eta}(\alpha) \geq 0$. By
    \Cref{theorem:artin_schreier_ramification}(1)--(2), the extension defined by
    $Z^p - Z = \alpha$ is trivial or unramified; either way, $\cP'^{[d]}$ ---
    and therefore $\cP^{[d]}$ --- is unramified at $\eta_\ndls$.
\end{proof}

We record the two cyclic conclusions obtained so far in the form used later.

\begin{cor}\label{cor:ramification_after_blowup_is_tame}
    In the general settings of a curve $C$, and $\cP \to U=C \setminus \mdls$ a $G$-torsor with ramification bounded by $\ndls = n P \subset \mdls$ at $P$, we have:
    \begin{enumerate}
        \item If $G=\Z/p\Z$ then $\cP^{[\deg \ndls]}$ is unramified at $\eta_\ndls$
        \item If $\cP$ is tamely ramified at $\ndls$ then $\cP^{[\deg \ndls]}$ is tamely ramified at $\eta_\ndls$
    \end{enumerate}
\end{cor}
\begin{proof}
    By \Cref{rem:algebraically_closed_harmless} we may assume $k$ is
    algebraically closed, so that $P \in C(k)$ and $\deg \ndls = n$. Part (1) is
    \Cref{theorem:wild_Zp_after_blowup}, and part (2) is
    \Cref{theorem:tame_after_blowup}.
\end{proof}

% ---------------------------------------------------------------------
\subsection{\texorpdfstring{The Wild Case $G = \Z/p^r\Z$}{The Wild Case G = Z/p\^{}rZ}}\label{subsection:wild_case_Zpr}
% ---------------------------------------------------------------------
% [Source: RamificationAfterBlowupPPower.tex, trimmed: the reduced-vector
%  material moved to 2.3.3, lemma:symmetric_regularity moved to the prelude;
%  internal references updated to the new per-case structure.]

Next, we generalize to the case $G = \Z/p^r\Z$.

\begin{theorem}\label{theorem:ramification_after_blowup_is_tame_for_p_power}
    Assume $G=\Z/p^r\Z$. Then $\cP^{[\deg (\ndls)]}$ is unramified at $\eta_{\ndls}$.
\end{theorem}

The proof occupies the remainder of this subsection. In contrast with what the phrase
``we generalize'' may suggest, the argument is \emph{not} an induction on $r$ with
\Cref{theorem:wild_Zp_after_blowup} as its base case: it is a direct
argument, treating all $r$ at once, in which each of the three steps of
\Cref{subsection:wild_case_Zp} --- realization of
the local torsor by an explicit equation over $\bG_m$, identification of the
symmetric-power cover at $\eta_\ndls$, and the degree estimate --- is replaced by its
Witt-vector analogue, using the theory collected in
\Cref{subsection:witt_vector_prerequisites} and
\Cref{subsection:asw_reduced_representatives}. In particular the case $r=1$ is
\emph{reproved} (as the special case of Witt vectors of length one) rather than used.
For an explanation of why the seemingly natural induction on $r$ fails --- in short:
the hypothesis bounds the ramification of $\cP$ in \emph{upper} numbering, while the
break of the top degree-$p$ layer $\cP \to \cP'$ is the largest \emph{lower}-numbering
break, which can be of size about $p^{r-1}d$ --- see \Cref{main:section:a_wrong_path}.

\subsubsection*{Standing reduction}

As in the preceding cases, we may assume $k$ is
algebraically closed (\Cref{rem:algebraically_closed_harmless}). We may further assume that $\cP \to U$ is \emph{totally
ramified} at $P$ with inertia equal to all of $G$: writing
$R = \widehat{\Oo_{C,P}}$, $K = \Frac(R)$, and letting
$H = \rho(G^0) \subseteq G$ be the image of the inertia subgroup under the
homomorphism $\rho \colon \Gal(K^{sep}/K) \to G$ classifying $\cP|_{\Spec K}$,
decompose $\cP \to \cP/H \to U$ as in \Cref{main:prop:quotient_torsors}, where
$\cP \to \cP/H$ is an $H$-torsor and $\cP/H \to U$ a $G/H$-torsor. Then
$\cP/H$ is unramified at $P$, hence extends over $P$; taking any point
$Q \in \cP/H$ over $P$ yields $\widehat{\Oo_{\cP/H, Q}} \cong \widehat{\Oo_{C,P}}$,
and since $k$ is algebraically closed the torsor $\cP/H$ is locally trivial at
$Q$. This replaces $(G, \cP \to U)$ by $(H, \cP \to \cP/H)$ with
$H \cong \Z/p^s\Z$ for some $s \leq r$; as the theorem is proved for all $r$
simultaneously, we may rename and assume $H = G$.

After this reduction, the restriction $\cQ := \cP_x$ of $\cP$ to the punctured formal
disc $\Spec k((x))$ at $P$ (with $x$ a uniformizer at $P$) is a \emph{connected}
$\Z/p^r\Z$-torsor: $\cQ = \Spec L$ for a cyclic, totally ramified extension $L/K$ of
degree $p^r$, with ramification bounded by $d := \deg \ndls$
(\Cref{definition:local_ramification_boundness}); equivalently, the largest
upper-numbering break $u$ of $L/K$ satisfies $u < d$ \emph{strictly}.

\subsubsection*{Step 1: realization over $\bG_m$}

The first step generalizes \Cref{lemma:realization_over_Gm_AS} from $\Z/p\Z$ to
$\Z/p^r\Z$, using the reduced polynomial representatives of
\Cref{subsection:asw_reduced_representatives}.

\begin{lemma}[Realization over $\bG_m$ for $W_r$]\label{lemma:realization_over_Gm_Wr}
    Let $k$ be algebraically closed of characteristic $p>0$ and let $\cQ$ be a
    \emph{connected} $\Z/p^r\Z$-torsor over $\Spec(k((x)))$ with ramification bounded
    by $d$ (\Cref{definition:local_ramification_boundness}). Then there is a reduced
    polynomial vector $\vec f = (f_0, \dots, f_{r-1})$, $f_j \in x^{-1}k[x^{-1}]$,
    such that, with $m_j := \deg_{x^{-1}} f_j$:
    \begin{enumerate}
        \item \label{item:realization_Wr_polebound} $p^{\,r-1-j}\, m_j < d$ for all
        $0 \leq j \leq r-1$; equivalently $m_j < d\, p^{\,j-(r-1)}$;
        \item \label{item:realization_Wr_torsor}
        $\cP' := \Spec\Bigl(k[x,x^{-1}][\vec X]\big/\bigl(\wp \vec X \ominus \vec f\bigr)\Bigr)$
        is a $\Z/p^r\Z$-torsor over $\bG_m$ with $\cP'_x \cong \cQ$;
        \item \label{item:realization_Wr_infty} $\cP'$ extends to a torsor over
        $\bP^1_k \setminus \{0\}$ --- in particular it is unramified at $\infty$ ---
        and has ramification bounded by $d$ at $0$.
    \end{enumerate}
\end{lemma}

\begin{proof}
    $\cQ = \Spec L$ with $L/K$ cyclic totally ramified of degree $p^r$ and largest
    upper break $u < d$ (\Cref{theorem:asw_theory}). Let
    $[\vec f_0] \in W_r(K)/\wp$ be its Artin--Schreier--Witt class and let $\vec f$ be
    a reduced polynomial representative
    (\Cref{prop:reduced_polynomial_representatives}).

    First, $f_0 \neq 0$: otherwise
    \Cref{cor:asw_class_order_and_first_component} would bound the order of the class
    by $p^{r-1}$, contradicting $[L:K] = p^r$. Hence $m_0 \geq 1$ and $p \nmid m_0$,
    and the Schmid--Witt break formula
    (\Cref{theorem:schmid_witt_break_formula}) applies:
    $\max_j p^{r-1-j} m_j = u < d$, which is \ref{item:realization_Wr_polebound}.

    For \ref{item:realization_Wr_torsor}: by \Cref{lemma:asw_etale_torsor} (with
    $A = k[x,x^{-1}] \ni f_j$), $\cP'$ is a $\Z/p^r\Z$-torsor over $\bG_m$; its
    restriction to $\Spec k((x))$ is the Artin--Schreier--Witt torsor of the class of
    $\vec f$, i.e.\ $\cQ$ (\Cref{theorem:asw_theory}).

    For \ref{item:realization_Wr_infty}: in the coordinate $y = x^{-1}$ at $\infty$ we
    have $f_j \in y\,k[y]$, so $\vec f \in W_r(k[y])$ and
    \Cref{lemma:asw_etale_torsor} (with $A = k[y]$) extends $\cP'$ to a torsor over
    $\A^1_k = \Spec k[y] = \bP^1_k \setminus \{0\}$; a torsor over the full local ring
    at $\infty$ is unramified there. The ramification bound at $0$ holds because
    $\cP'_x \cong \cQ$.
\end{proof}

\begin{rem*}\label{rem:realization_Wr_sanity}
    Connectedness of $\cQ$ is guaranteed in the application by the standing reduction
    (inertia equal to all of $G$). Note also the following sanity check: for every
    $j$ with $p^{\,r-1-j} \geq d$, the bound \ref{item:realization_Wr_polebound}
    forces $m_j = 0$, i.e.\ $f_j = 0$ --- all sufficiently low Witt components of a
    reduced representative vanish outright. This makes visible how strong the
    hypothesis $u < d$ is at the lower levels.
\end{rem*}

\subsubsection*{Step 2: identification of the symmetric-power cover}

We now generalize the invariant-ring computation of
\Cref{prop:as_cover_identification} to Witt coordinates. Fix
$\vec f \in W_r(R)$, $R = k[x, x^{-1}]$, and let
\[
S \;=\; R[\vec X]\big/\bigl(\wp \vec X \ominus \vec f\bigr)
\]
be the coordinate ring of the torsor $\cP' \to \bG_m$ of
\Cref{lemma:realization_over_Gm_Wr}. Write
$R^{\otimes_k d} = k[x_1^{\pm 1}, \dots, x_d^{\pm 1}]$ and
\[
S^{\otimes_k d}
= R^{\otimes_k d}[\vec X_1, \dots, \vec X_d]\big/
\bigl(\wp \vec X_1 \ominus \vec f(x_1), \dots, \wp \vec X_d \ominus \vec f(x_d)\bigr),
\]
where $\vec X_i = (X_{i,0}, \dots, X_{i,r-1})$ and $\vec f(x_i)$ is the image of
$\vec f$ under $x \mapsto x_i$. Exactly as in the scaffold of
\Cref{section:blowup_prelude}, the $d$-fold product torsor
$p_1^{-1}\cP' \otimes^G \cdots \otimes^G p_d^{-1}\cP'$ on $U^d$ is the quotient of
$\Spec S^{\otimes_k d}$ by $\KG \cong G^{d-1}$, embedded in $G^d$ as the kernel of the
sum map via $(g_1, \dots, g_{d-1}) \mapsto (g_1, g_2 - g_1, \dots, -g_{d-1})$, and
acting (in Witt coordinates, writing $\underline g$ for the image of $g \in G$ under
$\Z/p^r\Z \cong W_r(\F_p)$, \Cref{example:witt_of_Fp}) by
\[
(g_1, \dots, g_{d-1}) \colon \quad
\vec X_1 \mapsto \vec X_1 \oplus \underline g_1, \qquad
\vec X_i \mapsto \vec X_i \oplus \underline g_i \ominus \underline g_{i-1}
\ (1 < i < d), \qquad
\vec X_d \mapsto \vec X_d \ominus \underline g_{d-1}.
\]

\begin{prop}[identification of the product torsor, Witt case]\label{prop:witt_cover_identification}
    Set
    \[
    \vec Y := \bigoplus_{i=1}^{d} \vec X_i \in W_r\bigl(S^{\otimes_k d}\bigr),
    \qquad
    \vec \alpha := \bigoplus_{i=1}^{d} \vec f(x_i) \in W_r\bigl(R^{\otimes_k d}\bigr),
    \]
    where $\bigoplus$ denotes the $d$-fold Witt vector sum $\oplus$ (not a direct
    sum); in particular $\vec Y$ is a single Witt vector of length $r$. Then:
    \begin{enumerate}
        \item \label{item:witt_cover_invariance} Every component $Y_n$ of $\vec Y$ is
        $G^{d-1}$-invariant, and
        \[
        \bigl(S^{\otimes_k d}\bigr)^{G^{d-1}}
        \;\cong\;
        R^{\otimes_k d}[\vec Y]\big/\bigl(\wp \vec Y \ominus \vec \alpha\bigr)
        \]
        as $G$-torsors over $R^{\otimes_k d}$ (the residual $G = G^d/G^{d-1}$ acting
        by $\vec Y \mapsto \vec Y \oplus \underline g$).
        \item \label{item:witt_cover_symmetry} Every component $\alpha_n$ of
        $\vec \alpha$ is a symmetric Laurent polynomial, i.e.\ lies in
        $\bigl(R^{\otimes_k d}\bigr)^{S_d} = k[e_1, \dots, e_d, e_d^{-1}]$, and
        \[
        \Bigl(\bigl(S^{\otimes_k d}\bigr)^{G^{d-1}}\Bigr)^{S_d}
        \;\cong\;
        k[e_1, \dots, e_d, e_d^{-1}][\vec Y]\big/\bigl(\wp \vec Y \ominus \vec \alpha\bigr).
        \]
    \end{enumerate}
    Consequently, restricting to the complete local field
    $\widehat K_\eta = \kappa(\eta)((e_d))$,
    $\kappa(\eta) = k(e_1/e_d, \dots, e_{d-1}/e_d)$, of
    \Cref{eq:complete_field_at_generic_point},
    \[
    \cP'^{[d]}\big|_{\Spec \widehat K_\eta}
    \;\cong\;
    \Spec\; \widehat K_\eta[\vec Z]\big/\bigl(\wp \vec Z \ominus \vec \alpha\bigr).
    \]
\end{prop}

\begin{proof}
    \ref{item:witt_cover_invariance} An element
    $\sigma = (g_1, \dots, g_{d-1}) \in G^{d-1}$ acts as a ring automorphism of
    $S^{\otimes_k d}$: this is its \emph{original} action, defined on the generators
    $X_{i,n}$ by the twisted formulas above (which mix components through the carry
    polynomials of \Cref{eq:witt_addition_shape}). By functoriality of $W_r$
    (\Cref{subsection:witt_vector_prerequisites}), the induced map $W_r(\sigma)$ on
    $W_r(S^{\otimes_k d})$ is componentwise and commutes with $\oplus, \ominus$; on
    the generators it reproduces the twisted action,
    $W_r(\sigma)(\vec X_i) = \vec X_i \oplus \underline g_i \ominus \underline g_{i-1}$
    (with $\underline g_0 = \underline g_d = \vec 0$). Hence, using commutativity and
    associativity of $\oplus$, the Witt sum transforms by the telescoping constant
    \[
    W_r(\sigma)(\vec Y)
    \;=\; \bigoplus_{i=1}^{d} W_r(\sigma)(\vec X_i)
    \;=\; \vec Y \oplus \underline g_1 \oplus (\underline g_2 \ominus \underline g_1)
    \oplus \cdots \oplus (\ominus\, \underline g_{d-1})
    \;=\; \vec Y.
    \]
    Since $W_r(\sigma)$ is componentwise, this equality of $r$-tuples says exactly
    $\sigma(Y_n) = Y_n$ for every $n$, i.e.\ each $Y_n \in S^{\otimes_k d}$ is
    invariant. Since $\wp = F - \mathrm{id}$ is additive
    (\Cref{definition:asw_operator} ff.),
    \[
    \wp \vec Y = \bigoplus_{i=1}^{d} \wp \vec X_i
    = \bigoplus_{i=1}^{d} \vec f(x_i) = \vec \alpha .
    \]
    Hence there is a well-defined $R^{\otimes_k d}$-algebra homomorphism
    \[
    \phi \colon\; T := R^{\otimes_k d}[\vec Z]\big/\bigl(\wp \vec Z \ominus \vec \alpha\bigr)
    \;\longrightarrow\; \bigl(S^{\otimes_k d}\bigr)^{G^{d-1}},
    \qquad \vec Z \mapsto \vec Y .
    \]
    By \Cref{lemma:asw_etale_torsor}, $\Spec T$ is a $G$-torsor over $U^d$. The
    target is the product torsor (the contracted product of the $p_i^{-1}\cP'$ along
    the sum map $G^d \to G$), a $G$-torsor over $U^d$ by
    \Cref{main:prop:symmetric_power_torsor} and the surrounding discussion; its residual
    $G$-action is induced by $(g, 0, \dots, 0) \in G^d$, i.e.\
    $\vec X_1 \mapsto \vec X_1 \oplus \underline g$, under which
    $\vec Y \mapsto \vec Y \oplus \underline g$. So $\Spec \phi$ is a $G$-equivariant
    morphism of $G$-torsors over $U^d$, hence an isomorphism by
    \Cref{main:prop:torsor_morphism_iso}.

    \ref{item:witt_cover_symmetry} $S_d$ acts on $S^{\otimes_k d}$ by permuting the
    tensor factors, i.e.\ $x_i \mapsto x_{\sigma(i)}$,
    $\vec X_i \mapsto \vec X_{\sigma(i)}$. As in
    \ref{item:witt_cover_invariance}, each $\sigma \in S_d$ is a ring automorphism,
    so $W_r(\sigma)$ is componentwise and commutes with $\oplus$. Since $\oplus$ is
    commutative and associative, permuting the summands of a Witt sum leaves it
    unchanged; thus $\vec Y$ is $S_d$-invariant (componentwise, as before) and
    $W_r(\sigma)(\vec \alpha) = \bigoplus_i \vec f(x_{\sigma(i)}) = \vec \alpha$, so
    each $\alpha_n \in R^{\otimes_k d}$ is a symmetric Laurent polynomial, hence lies
    in $k[e_1, \dots, e_d, e_d^{-1}]$ (the $S_d$-invariants of $R^{\otimes_k d}$, as
    recorded in the scaffold). By
    \Cref{lemma:asw_etale_torsor}, $T$ is free over $R^{\otimes_k d}$ with basis the
    monomials $\prod_n Y_n^{a_n}$, $0 \leq a_n < p$, each of which is
    $S_d$-invariant; writing an element in this basis, it is $S_d$-invariant if and
    only if all its coefficients are. Therefore
    \[
    T^{S_d} = \bigl(R^{\otimes_k d}\bigr)^{S_d}[\vec Y]\big/\bigl(\wp \vec Y \ominus \vec \alpha\bigr)
    = k[e_1, \dots, e_d, e_d^{-1}][\vec Y]\big/\bigl(\wp \vec Y \ominus \vec \alpha\bigr).
    \]

    The final statement follows exactly as in the $r=1$ case: by definition
    $\cP'^{[d]}$ is the quotient by $\Xi = \KG \rtimes S_d$ computed above
    (\Cref{main:cor:sym-power-quotient}), and its restriction to $\Spec \widehat K_\eta$
    along $k[e_1, \dots, e_d, e_d^{-1}] \to \widehat K_\eta$ is the
    Artin--Schreier--Witt cover of the image of $\vec \alpha$ in
    $W_r(\widehat K_\eta)$.
\end{proof}

\begin{rem*}
    A pure restriction argument on $H^1$'s (``the
    class of the product torsor is the sum of the pulled-back characters'')
    identifies the cover over $k(x_1, \dots, x_d)$ but does \emph{not} by itself
    descend it to the symmetric base $k(e_1, \dots, e_d)$; the invariant-theoretic
    computation above --- like its Kummer and Artin--Schreier counterparts,
    \Cref{prop:kummer_cover_identification,prop:as_cover_identification} --- is
    what effects the descent.
\end{rem*}

\subsubsection*{Step 3: the degree count}

Throughout this paragraph, $\vec f$ is the reduced polynomial vector of
\Cref{lemma:realization_over_Gm_Wr}, so $m_j = \deg_{x^{-1}} f_j$ satisfies
$m_j < d\, p^{\,j-(r-1)}$, and $y_i := 1/x_i$, so that $f_j(x_i) \in y_i\, k[y_i]$ is
a polynomial in $y_i$ of degree $m_j$.

\begin{theorem}[regularity of all Witt components]\label{theorem:witt_components_regularity}
    With $\vec f$ as in \Cref{lemma:realization_over_Gm_Wr}, every component
    $\alpha_n$ $(0 \leq n \leq r-1)$ of
    $\vec \alpha = \bigoplus_{i=1}^{d} \vec f(x_i)$ lies in $\kappa(\eta)$; in
    particular $\vec \alpha \in W_r\bigl(\widehat\Oo_\eta\bigr)$.
\end{theorem}

\begin{proof}
    Fix $n \leq r-1$ and a monomial $\prod_{i,j} \bigl(f_j(x_i)\bigr)^{e_{ij}}$ of
    $\alpha_n$, viewed as a universal isobaric polynomial in the $f_j(x_i)$ of
    weight $p^n$ (\Cref{lemma:witt_addition_isobaric}, with $f_j(x_i)$ of weight
    $p^j$); set $E_j := \sum_i e_{ij}$, so $\sum_j E_j p^{\,j} = p^{\,n}$. Expanding
    each $f_j(x_i)$ as a polynomial in $y_i$ of degree $m_j$, every resulting
    monomial in the $y_i$ has total degree
    \[
    \leq\; \sum_{i,j} e_{ij}\, m_j \;=\; \sum_j E_j\, m_j
    \;<\; \sum_j E_j\, d\, p^{\,j-(r-1)}
    \;=\; d\, p^{-(r-1)} \sum_j E_j p^{\,j}
    \;=\; d\, p^{\,n-(r-1)} \;\leq\; d,
    \]
    strictly: some $E_j > 0$ (as $\sum_j E_j p^j = p^n > 0$), and the pole bound
    $m_j < d\, p^{\,j-(r-1)}$ of
    \Cref{lemma:realization_over_Gm_Wr}\ref{item:realization_Wr_polebound} is strict
    even when $m_j = 0$. Thus every monomial of $\alpha_n$ --- a symmetric polynomial
    in $y_1, \dots, y_d$ by
    \Cref{prop:witt_cover_identification}\ref{item:witt_cover_symmetry} --- has total
    $y$-degree $< d$, and \Cref{lemma:symmetric_regularity} gives
    $\alpha_n \in \kappa(\eta)$.
\end{proof}

\begin{rem*}
    The bound is exact: the weight $p^{\,j-(r-1)}$ is paired against the isobaricity
    constraint $\sum_j E_j p^{\,j} = p^{\,n}$, so the estimate
    $d\, p^{\,n-(r-1)} \leq d$ is independent of how the Witt carries distribute the
    exponents $E_j$; the hypothesis ``break $< d$ at $P$'' is used at \emph{every}
    Witt level at once. The sanity check of
    \Cref{lemma:realization_over_Gm_Wr} records the extreme case: components with
    $p^{\,r-1-j} \geq d$ vanish outright in the reduced representative.
\end{rem*}

\subsubsection*{Proof of the theorem}

\begin{proof}[Proof of \Cref{theorem:ramification_after_blowup_is_tame_for_p_power}]
    By the standing reduction above we may assume $k$ algebraically closed and
    $\cP \to U$ totally ramified at $P$ with inertia all of $G$, so that
    $\cQ := \cP_x$ is a connected $\Z/p^r\Z$-torsor over $\Spec k((x))$ with
    ramification bounded by $d = \deg \ndls$.

    \Cref{lemma:realization_over_Gm_Wr} produces a torsor $\cP' \to \bG_m$, defined
    by a reduced polynomial vector $\vec f$ with $p^{r-1-j} m_j < d$, with
    $\cP'_x \cong \cQ$. By \Cref{cor:reduction_to_P1} (valid for arbitrary finite
    abelian $G$), the ramification of $\cP^{[d]}$ at $\eta_\ndls$ agrees with that of
    $\cP'^{[d]}$ at $\eta_{d \cdot 0}$; so we may assume $C = \bP^1_k$,
    $\mdls = d \cdot 0 + \infty$, $\ndls = d \cdot 0$, $\cP = \cP'$.

    By \Cref{prop:witt_cover_identification},
    \[
    \cP'^{[d]}\big|_{\Spec \widehat K_\eta}
    \;\cong\;
    \Spec\; \widehat K_\eta[\vec Z]\big/\bigl(\wp \vec Z \ominus \vec \alpha\bigr),
    \qquad
    \vec \alpha = \bigoplus_{i=1}^{d} \vec f(x_i),
    \]
    with $\widehat K_\eta = \kappa(\eta)((e_d))$ the complete local field at
    $\eta_\ndls$ (\Cref{eq:complete_field_at_generic_point}) and
    $\widehat\Oo_\eta = \kappa(\eta)[[e_d]]$ its valuation ring.

    By \Cref{theorem:witt_components_regularity},
    $\vec \alpha \in W_r(\kappa(\eta)) \subset W_r(\widehat\Oo_\eta)$. Hence, by
    \Cref{lemma:asw_etale_torsor} applied over $A = \widehat\Oo_\eta$, the
    Artin--Schreier--Witt cover $\wp \vec Z = \vec \alpha$ is finite \'etale over
    $\widehat\Oo_\eta$, and $\cP'^{[d]}\big|_{\Spec \widehat K_\eta}$ is its generic
    fibre: every field factor of
    $\widehat K_\eta[\vec Z]/(\wp \vec Z \ominus \vec \alpha)$ is the fraction field
    of a finite \'etale $\widehat\Oo_\eta$-algebra, i.e.\ an unramified extension of
    $\widehat K_\eta$ with respect to the discrete valuation ring
    $\widehat\Oo_\eta$. By \Cref{definition:tamely_ramified_in_codim_1},
    $\cP^{[d]}$ is unramified at $\eta_\ndls$.
\end{proof}

\begin{rem*}
    The argument proves all $r$ simultaneously; it reproves the case $r=1$
    (\Cref{theorem:wild_Zp_after_blowup}) as the special case of a
    vector of length one and never invokes it.
\end{rem*}

% ---------------------------------------------------------------------
\subsection{The General Finite Abelian Case}\label{subsection:general_abelian_case}
% ---------------------------------------------------------------------
% [NEW: this subsection replaces the one-line stub "Finally, Proof of Theorem
%  ...: Idea: decompose G by the structure theorem, and finish by the above."]

We now assemble the general case from the three cyclic computations. The
mechanism is the structure theorem for finite abelian groups together with the
compatibility of the symmetric power with fiber products
(\Cref{lemma:sym_power_product}).

\begin{proof}[Proof of \Cref{theorem:torsors_after_blowup_is_tame}]
    Fix a single-point submodulus $\ndls = n_i P_i$ of $\mdls$, write $P = P_i$
    and $d = d_\ndls$. By \Cref{rem:algebraically_closed_harmless} we may assume
    $k$ is algebraically closed, so $d = n_i$.

    By the structure theorem for finite abelian groups, fix an isomorphism
    \[
        G \;\cong\; \prod_{j=1}^{s} \Z/q_j\Z,
        \qquad q_j = \ell_j^{\,r_j} \ \text{prime powers},
    \]
    and let $\pi_j : G \to \Z/q_j\Z$ be the projections. Let
    $\cP_j := \cP / \ker(\pi_j)$ be the quotient torsor
    (\Cref{main:prop:quotient_torsors}); it is a $\Z/q_j\Z$-torsor over $U$, and at
    each point of $\mathrm{Supp}(\mdls)$ its classifying character is
    $\pi_j \circ \rho$, where $\rho$ is that of $\cP$.

    \emph{Step 1: $\cP$ is the fiber product of its primary components.} The
    quotient maps $\cP \to \cP_j$ combine to a morphism
    \[
        \cP \longrightarrow \cP_1 \times_U \cP_2 \times_U \cdots \times_U \cP_s
    \]
    which is equivariant for the isomorphism $G \cong \prod_j \Z/q_j\Z$. The
    target, being a fiber product of $\Z/q_j\Z$-torsors, is a
    $\prod_j \Z/q_j\Z$-torsor, so the morphism is an isomorphism of $G$-torsors
    (\Cref{main:prop:torsor_morphism_iso}).

    \emph{Step 2: each component satisfies the hypothesis of one of the three
    cyclic cases.} Since $\rho(G^{\,n_i}_{\widehat{K}_{P_i}}) = \{1\}$ implies
    $(\pi_j \circ \rho)(G^{\,n_i}_{\widehat{K}_{P_i}}) = \{1\}$, each $\cP_j$ has
    ramification bounded by $\mdls$; in particular, bounded by $\ndls = dP$ at
    $P$. Moreover:
    \begin{itemize}
        \item If $\ell_j \neq p$: the image of the wild inertia subgroup under
        $\pi_j \circ \rho$ is a $p$-group inside $\Z/q_j\Z$, whose order is prime
        to $p$; so it is trivial, and $\cP_j$ is \emph{tamely ramified} at every
        point of $\mathrm{Supp}(\mdls)$ --- in particular at $P$.
        \item If $\ell_j = p$: $\cP_j$ is a $\Z/p^{r_j}\Z$-torsor with
        ramification bounded by $\ndls$ at $P$.
    \end{itemize}

    \emph{Step 3: the cyclic cases.} By \Cref{theorem:tame_after_blowup}, for
    every $j$ with $\ell_j \neq p$ the torsor $\cP_j^{[d]}$ is tamely ramified at
    $\eta_\ndls$. By \Cref{theorem:ramification_after_blowup_is_tame_for_p_power},
    for every $j$ with $\ell_j = p$ the torsor $\cP_j^{[d]}$ is unramified ---
    in particular tamely ramified --- at $\eta_\ndls$.

    \emph{Step 4: reassembly.} Applying \Cref{lemma:sym_power_product}
    repeatedly to Step 1,
    \[
        \cP^{[d]}
        \;\cong\;
        \cP_1^{[d]} \times_{U^{(d)}} \cP_2^{[d]} \times_{U^{(d)}} \cdots
        \times_{U^{(d)}} \cP_s^{[d]}.
    \]
    By \Cref{lemma:ramification_of_fiber_product} (applied repeatedly at the
    prime divisor $(E_\ndls, \eta_\ndls)$ of $X_\ndls$), a fiber product of
    tamely ramified torsors is tamely ramified. Hence $\cP^{[d]}$ is tamely
    ramified at $\eta_\ndls$, as claimed.
\end{proof}

\begin{rem*}
    The proof shows slightly more: the wild part of $\cP$ contributes
    \emph{unramified} factors at $\eta_\ndls$, so the ramification of
    $\cP^{[d]}$ at $\eta_\ndls$ is entirely accounted for by the tame,
    prime-to-$p$ part of $\cP$ --- with ramification index dividing the order
    of the tame quotient of the local character at $P$, as computed in
    \Cref{theorem:tame_after_blowup}.
\end{rem*}

% ---------------------------------------------------------------------------
%  2.6  Ramification on Products of Blowups
%  [Source: RamificationAfterBlowupIsTame.tex lines 837--957, unchanged except
%   that the duplicated "Affine Case" computation is replaced by a reference
%   to prop:blowup_local_ring_at_exceptional in the prelude.]
% ---------------------------------------------------------------------------

\section{Ramification on Products of Blowups}\label{subsection:ram_prod_analysis}
We conclude this chapter by proving some auxiliary results that will be useful in \Cref{main:section:gcft}.

Let $X$ and $Y$ be smooth schemes over a field $k$, and let $x \in X$ and $y \in Y$ be closed points.
We denote the blowups of these schemes at the given points by
$\pi_X: \Bl_x(X) \to X$ and $\pi_Y: \Bl_y(Y) \to Y$. Furthermore, let $\pi_{X \times Y}: \Bl_{(x,y)}(X \times_k Y) \to X \times_k Y$
be the blowup of the product scheme at the point $(x,y)$.
We denote by $E_X, E_Y$, and $E_{X \times Y}$ the respective exceptional divisors, and let $\eta_X, \eta_Y$, and $\eta_{X \times Y}$ be their generic points.

In this section, we establish the following result concerning the stability of ramification bounds under the external product of torsors.

\begin{prop}\label{prop:ramification-external-product}
Let $G$ be a finite abelian group. Suppose $\cG_X$ and $\cG_Y$ are $G$-torsors defined on open subsets $U_X \subset \Bl_x(X)$ and $U_Y \subset \Bl_y(Y)$ that are disjoint
from the exceptional divisors.
If the ramification of $\cG_X$ at $\eta_X$ and $\cG_Y$ at $\eta_Y$ is bounded by $r$, then the external product torsor
\[
\cG_{X \times Y} := \text{pr}_1^{-1} \cG_X \otimes \text{pr}_2^{-1} \cG_Y
\]
has ramification bounded by $r$ at the generic point $\eta_{X\times Y}$ of the exceptional divisor in the product blowup.
\end{prop}

The proposition is purely local in nature; it suffices to consider the case where $X$ and $Y$ are affine.
More precisely, by the smoothness of $X$ and $Y$, we may restrict our attention to open neighborhoods of $x$ and $y$ that are étale over affine spaces.
The rest of this section treats that case. So let $X = \mathbb{A}^n_k$ and
$Y = \mathbb{A}^m_k$ with origins $x = 0$ and $y = 0$; the local rings, residue
fields and valuations at the generic points $\eta_X$ and $\eta_Y$ of the
exceptional divisors $E_X \subset \Bl_0(X)$ and $E_Y \subset \Bl_0(Y)$ are the
ones computed in \Cref{prop:blowup_local_ring_at_exceptional}.

\subsection*{The Product Case}
As before, $\text{Bl}_0(X) \subset X \times \mathbb{P}^{n-1}$ and $\text{Bl}_0(Y) \subset Y \times \mathbb{P}^{m-1}$ have exceptional divisors $E_X$ and $E_Y$ respectively.
Consider the product $X \times_k Y \cong \mathbb{A}_k^{n+m}$.
The blowup of the product at the origin $(0,0)$, denoted $\text{Bl}_{(0,0)}(X \times_k Y)$, is a subscheme of $(X \times Y) \times \mathbb{P}^{n+m-1}$ defined by:

$$\begin{cases}
x_i w_j = x_j w_i & 1 \le i, j \le n \\
y_k w_{n+l} = y_l w_{n+k} & 1 \le k, l \le m \\
x_i w_{n+j} = y_j w_i & 1 \le i \le n, \,\, 1 \le j \le m
\end{cases}$$
where $[w_1 : \dots : w_{n+m}]$ are the homogeneous coordinates of $\mathbb{P}^{n+m-1}$. The exceptional divisor $E_{X \times Y}$ is isomorphic to $\mathbb{P}^{n+m-1}$.

\subsection*{Comparison of Blowups}

Both $\text{Bl}_0(X) \times_k \text{Bl}_0(Y)$ and $\text{Bl}_{(0,0)}(X \times_k Y)$ are birational to $X \times_k Y$.
They share a common dense open set $\tilde{U}$ defined by the condition that neither the $X$-coordinates nor the $Y$-coordinates vanish simultaneously in the projective space:
$$\tilde{U} = \{ ((x,y), [w_1: \dots : w_{n+m}]) \mid (w_1, \dots, w_n) \neq 0 \text{ and } (w_{n+1}, \dots, w_{n+m}) \neq 0 \}$$This yields a diagram of open immersions:
$$
   \begin{tikzcd}
        & \tilde{U} \arrow[dl, "f_1"'] \arrow[dr, "f_2", hook] & \\
        \text{Bl}_0(X) \times_k \text{Bl}_0(Y) & & \text{Bl}_{(0,0)}(X \times_k Y)
    \end{tikzcd}
$$

    where $f_1$ maps the coordinates to the respective projectivizations $[w_1: \dots : w_n]$ and $[w_{n+1}: \dots : w_{n+m}]$.
    Also note that the generic point $\eta_{X \times Y}$  of $E_{X \times_k Y}$ is inside $\tilde{U}$.

\subsection*{Extensions of DVRs}
Let $S$ be the local ring of the generic point $\eta_{X\times Y}$ of $E_{X \times Y}$ in $\tilde{U}$.
On the chart where $w_1 \neq 0$ and $w_{n+1} \neq 0$, we have
$x_1 = (\frac{w_1}{w_{n+1}}) y_1$.
Since $\frac{w_1}{w_{n+1}}$ is a unit in this chart, $x_1$ and $y_1$ are equivalent as uniformizers.
We have:
$$
S=k\left[x_1, \frac{w_2}{w_1} \cdots, \frac{w_n}{w_1}, \frac{w_{n+1}}{w_1} \cdots \frac{w_{n+m}}{w_1}\right]_{(x_1)} =
k\left[\frac{w_1}{w_{n+1}}, \frac{w_2}{w_{n+1}} \cdots \frac{w_n}{w_{n+1}}, y_1 \cdots \frac{w_{n+m}}{w_{n+1}}\right]_{(y_1)}
$$

$$k(\eta_{X \times Y}) = k\left(\frac{w_2}{w_1} \cdots, \frac{w_n}{w_1}, \frac{w_{n+1}}{w_1} \cdots \frac{w_{n+m}}{w_1}\right)$$

Let $R_X$ be the local ring of the exceptional divisor $E_X$ in $\text{Bl}_0(X)$
(computed in \Cref{prop:blowup_local_ring_at_exceptional}).
The pullback of $E_X \times_k \text{Bl}_0(Y)$ along $f_1$ induces an extension of DVRs $R_X \hookrightarrow S$. Which is:
\begin{enumerate}
    \item Weakly Unramified: $x_1$ is a uniformizer in both $R_X$ and $S$, so the ramification index is $e=1$.
    \item Residually Transcendental: The residue field extension $\kappa(\eta_X) \subset \kappa(\eta)$ is:
    $$k\left(\frac{w_2}{w_1}, \dots, \frac{w_n}{w_1}\right) \subset k\left(\frac{w_2}{w_1}, \dots, \frac{w_n}{w_1}, \frac{w_{n+1}}{w_1}, \dots, \frac{w_{n+m}}{w_1}\right)$$
    Hence separable.
\end{enumerate}
Since this extension is generated by transcendental elements, it is separable and formally smooth at the maximal ideal (\stackstag{09E7}).

\subsection*{Ramification of $G$-Torsors on the Product}
Let $G$ be a finite abelian group.
Let $\mathcal{Q}$ be a $G$-torsor on an open $U \subset \text{Bl}_0(X)$ disjoint from $E_X$.
Let $V \subset \text{Bl}_0(Y)$ be an open subscheme, and let $\pi_X: \text{Bl}_0(X) \times_k V \to \text{Bl}_0(X)$
be the projection onto the first factor.
By restricting this projection to $U \times_k V$, we obtain the pullback $G$-torsor:
$$\pi_X^{-1}(\mathcal{Q}) \cong \mathcal{Q} \times_k V$$
which is defined on the open subset $U \times_k V \subset \text{Bl}_0(X) \times_k \text{Bl}_0(Y)$.
The extension of local rings $R_X \to S$ is weakly unramified (the ramification index $e=1$) and residually transcendental with separable residue field extension.
Under these conditions the ramification filtration is preserved.
Therefore, the pullback $\pi_X^{-1}(\mathcal{Q})$ has ramification bounded by $r$ at the generic point $\eta_{X \times Y}$ of the exceptional divisor $E_{X \times Y}$ if and only if
the original torsor $\mathcal{Q}$ has ramification bounded by $r$ at the generic point $\eta_X$ of $E_X$.

And we finish by \Cref{lemma:bounding_ramification_of_contracted_product}.
